\subsection{Proper partial maps}

In this section, X and Y are locally compact topological spaces. We denote by \(X'\) (resp. \(Y'\) ) the compact space obtained from X (resp. Y) by adjoining a point at infinity \(\omega_{X}\) (resp. \(\omega_{Y}\) ) (TG, I, p.~67--68). We identify \(X'\) and \(Y'\) with \(\mathsf{X}'(\mathcal{C}_{0}(\mathsf{X}))\) and \(\mathsf{X}'(\mathcal{C}_{0}(\mathsf{Y}))\) respectively (Corollary 2 of I, p.~32).

DEFINITION 1. --- A proper partial map from X to Y is a correspondence \(f = (\Gamma, \mathrm{X}, \mathrm{Y})\) (E, II, p.~10, def. 2) between X and Y such that

\begin{enumerate}
\def\labelenumi{(\roman{enumi})}
\item
  The graph \(\Gamma\) is functional;
\item
  The domain of definition of f is an open set U of X;
\item
  The map \(x\mapsto f(x)\) from U to Y is proper.
\end{enumerate}

The identity map on X is a proper partial map from X into X. Let Z be a locally compact topological space and \(f\) (resp. \(g\) ) a proper partial map from X to Y (resp. from Y to Z). Then the composite correspondence \(g \circ f\) (E, II, p.~11, def. 6) is a proper partial map from X into Z (TG, I, p.~72, prop. 3, and p.~73, prop. 5).

Lemma 1. --- For every proper partial map f from X to Y, with domain of definition U, denote by \(\widetilde{f}\) the mapping from X' to Y' defined by \(\widetilde{f}(x) = f(x)\) if \(x \in \mathrm{U}\) and \(\widetilde{f}(x) = \omega_{\mathrm{Y}}\) if \(x \notin \mathrm{U}\) ; it is continuous.

The map \(f \mapsto \widetilde{f}\) is a bijection between the set of proper partial maps \(f\) from X to Y and the set of continuous maps \(g\) from X' to Y' such that \(g(\omega_{\mathrm{X}}) = \omega_{\mathrm{Y}}\) .

Let \(f\) a proper partial map from X to Y, and let U be its domain. Let us prove that the map \(\widetilde{f}\) is continuous. It is continuous at every point of U, since U is open in \(\mathrm{X}'\) Let us prove that it is also continuous at every point \(x\) of \(\mathrm{X}' - \mathrm{U}\) ; we then have \(\widetilde{f}(x) = \omega_{\mathrm{Y}}\) Let V be an open neighborhood of \(\omega_{\mathrm{Y}}\) into \(\mathrm{Y}'\) ; let us prove that \(\widetilde{f}^{-1}(\mathrm{V})\) is a neighborhood of \(x\) By definition of the topological space \(\mathrm{Y}'\) , one may assume that V is of the form \(Y^{\prime}-K\) where K is a compact subset of Y. Since f defines a proper map from U to Y, the set \(f^{-1}(K)\) is compact in U (TG, I, p.~77, prop. 6), hence in \(X^{\prime}\) This is in particular a closed subset of \(X^{\prime}\) and \(\tilde{f}^{-1}(V)=X^{\prime}-f^{-1}(K)\) is an open subset of \(X^{\prime}\) and is therefore a neighborhood of x.

Conversely, let \(g \colon X' \to Y'\) be a continuous map such that \(g(\omega_{\mathrm{X}}) = \omega_{\mathrm{Y}}\) and \(\Gamma_g \subset X' \times Y'\) its graph. The set \(U = X - \overline{g}^{-1}(\omega_{\mathrm{Y}})\) is open in \(X\) . The correspondence \(f = (\Gamma_g \cap (U \times Y), X, Y)\) is a proper partial map from \(X\) into \(Y\) such that \(\widetilde{f} = g\) , and it is the only one.

We will identify the proper partial maps from X to Y with the continuous maps from \(X'\) into \(Y'\) that map \(\omega_{X}\) to \(\omega_{Y}\) . In particular, the proper maps from X to Y are the proper partial maps with domain X; they are identified with the continuous maps f from \(X'\) into \(Y'\) such that \(f^{-1}(\omega_{\mathrm{Y}}) = \{\omega_{\mathrm{X}}\}\) . If X is compact, these are simply the continuous maps from X to Y.

Let A be a commutative complex Banach algebra. Recall that \(X'(A)\) is identified with the compact space obtained from \(X(A)\) by adjoining a point at infinity (I, p.~29, corollary).

PROPOSITION 2. --- Let A and B be commutative complex Banach algebras. For every algebra homomorphism \(\pi\colon\) A \(\to\) B, the map \(\mathsf{X}'(\pi)\) is a proper partial map from \(\mathsf{X}(\mathsf{B})\) into \(\mathsf{X}(\mathsf{A})\) .

Indeed, \(\mathsf{X}'(\pi)\) is a continuous map from \(\mathsf{X}'(\mathrm{B})\) into \(\mathsf{X}'(\mathrm{A})\) (I, p.~10). The point at infinity of \(\mathsf{X}'(\mathrm{B})\) (resp. of \(\mathsf{X}'(\mathrm{A})\) ) is the zero character, and we have \(\mathsf{X}'(\pi)(0)=0\) .

PROPOSITION 3. --- a) For every proper partial map \(\varphi\) from X to Y, the map \(f \mapsto f \circ \varphi\) of \(\mathcal{C}(\mathrm{Y}')\) into \(\mathcal{C}(\mathrm{X}')\) induces an algebra morphism \(\varphi^*\) of \(\mathcal{C}_0(\mathrm{Y})\) into \(\mathcal{C}_0(\mathrm{X})\) ;

\begin{enumerate}
\def\labelenumi{\alph{enumi})}
\setcounter{enumi}{1}
\tightlist
\item
  The map \(\varphi \mapsto \varphi^{*}\) is a bijection from the set of proper partial maps from X to Y onto the set of algebra morphisms from \(\mathcal{C}_0(\mathrm{Y})\) into \(\mathcal{C}_0(\mathrm{X})\) . Its inverse bijection is the map \(\pi \mapsto \mathsf{X}'(\pi)\) .
\end{enumerate}

Let us prove a). Let \(\varphi\) be a proper partial map from X to Y, identified with a continuous map from X' to Y' such that \(\varphi(\omega_{\mathrm{X}}) = \omega_{\mathrm{Y}}\) . For \(f \in \mathcal{C}_0(\mathrm{Y})\) , one has \((f \circ \varphi)(\omega_{\mathrm{X}}) = f(\omega_{\mathrm{Y}}) = 0\) , hence the map \(\varphi^*\) is well defined. It is an algebra morphism.

Let us prove that \(\mathsf{X}'(\varphi^{*})\) is identified with \(\varphi\) . Let \(x \in X\) . For every function \(f \in \mathcal{C}_{0}(\mathrm{Y})\) , the character \(\mathsf{X}'(\varphi^{*})(\mathrm{ev}_{x})\) associates to f the complex number

\[
\left(\operatorname{ev} _ {x} \circ \varphi^ {*}\right) (f) = \operatorname{ev} _ {x} (f \circ \varphi) = f (\varphi (x)),
\]

hence \(\mathsf{X}'(\varphi^{*})(\mathrm{ev}_{x}) = \mathrm{ev}_{\varphi(x)}\) . This proves the assertion.

Conversely, let \(\pi\colon\mathcal{C}_{0}(\mathrm{Y})\to\mathcal{C}_{0}(\mathrm{X})\) be an algebra morphism. Let us prove that \(\mathsf{X}^{\prime}(\pi)^{*}=\pi\) . Let \(f\in\mathcal{C}_{0}(\mathrm{Y})\) , and denote \(g=\mathsf{X}^{\prime}(\pi)^{*}(f)\in\mathcal{C}_{0}(\mathrm{X})\) . For every \(x\in\mathsf{X}\) , one has

\[
g (x) = \left(f \circ \mathrm{X} ^ {\prime} (\pi)\right) (x) = \operatorname{ev} _ {\mathrm{X} ^ {\prime} (\pi) (x)} (f) = \left(\operatorname{ev} _ {x} \circ \pi\right) (f) = \pi (f) (x),
\]

since \(\mathsf{X}^{\prime}(\pi)\) satisfies \(\mathrm{ev}_{\mathsf{X}^{\prime}(\pi)(x)} = \mathrm{ev}_{x} \circ \pi\) . We therefore have \(g = \pi(f)\) , which allows us to conclude that \(\mathsf{X}^{\prime}(\pi)^{*} = \pi\) .

In a completely similar way, we have:

PROPOSITION 4. --- Suppose that X and Y are compact. Identify the space X (resp. the space Y) with X( \(\mathcal{C}(X)\) ) (resp. X( \(\mathcal{C}(Y)\) )) (corollary 3 of I, p.~33).

\begin{enumerate}
\def\labelenumi{\alph{enumi})}
\item
  For every continuous map \(\varphi\colon\mathrm{X}\to\mathrm{Y}\) , the map \(\varphi^{*}\colon f\mapsto f\circ\varphi\) is an algebra morphism from \(\mathcal{C}(\mathrm{Y})\) into \(\mathcal{C}(\mathrm{X})\) ;
\item
  The maps \(\varphi \mapsto \varphi^{*}\) and \(\pi \mapsto X(\pi)\) are reciprocal bijections between the set of continuous maps from \(X\) into \(Y\) and the set of algebra morphisms from \(\mathcal{C}(Y)\) into \(\mathcal{C}(X)\) .
\end{enumerate}

Remark. --- *In the language of category theory, the preceding results are interpreted as follows. Let G be the category whose objects are locally compact topological spaces and whose morphisms are proper partial maps. The functor X ↦ \(\mathcal{C}_0(X)\) is a contravariant, fully faithful functor from the category G to the category of commutative complex Banach algebras. Moreover, A ↦ X(A) is a contravariant functor from the category of commutative complex Banach algebras to the category G. If to a locally compact topological space X one associates the homeomorphism

\[
\operatorname{ev} \colon X \to X (\mathscr {C} _ {0} (X)),
\]

one obtains an isomorphism from the identity functor of the category G to the composite functor \(X \mapsto X(\mathcal{C}_{0}(X))\) .

It is not true that the composite functor \(A \mapsto \mathcal{C}_{0}(X(A))\) is isomorphic to the identity functor of the category of commutative complex Banach algebras (cf.~example 2 of I, p.~36 and exercise 2 of I, p.~155). One will nevertheless see a statement of this type for commutative stellar algebras (number 5 of I, p.~107).*

\subsection{Gelfand transformation}

Let A be a commutative Banach algebra. Recall that, for every \(x \in A\) , we denote by \(\mathcal{G}_{\mathrm{A}}(x)\) , or \(\mathcal{G}(x)\) , the function \(\chi \mapsto \chi(x)\) on \(\mathsf{X}(\mathsf{A})\) , that \(\mathcal{G}(x)\) is called the Gelfand transform of x, and that the map \(x \mapsto \mathcal{G}(x)\) is called the Gelfand transformation (cf.~def. 5 of I, p.~7). One therefore has by definition:

\[
\mathcal {G} (x) (\chi) = \chi (x).
\]

Examples. --- 1) Let X be a locally compact topological space and consider the commutative Banach algebra \(\mathcal{C}_{0}(\mathrm{X})\) (example 3 of I, p.~17 and number 2 of I, p.~31). By cor. 1 of I, p.~32, the space of characters \(\mathsf{X}(\mathcal{C}_{0}(\mathbf{X}))\) is identified with X by means of the map associating to an element \(x \in \mathbf{X}\) the character \(f \mapsto f(x)\) of \(\mathcal{C}_{0}(\mathbf{X})\) , and the Gelfand transformation of \(\mathcal{C}_{0}(\mathbf{X})\) is then identified with the identity map.

\begin{enumerate}
\def\labelenumi{\arabic{enumi})}
\setcounter{enumi}{1}
\tightlist
\item
  Let \(n \geqslant 0\) be an integer. Let \(A_n\) be the algebra of functions \(f: [0,1] \to K\) admitting continuous derivatives in \([0,1]\) up to order \(n\) . Equipped with the norm
\end{enumerate}

\[
\| f \| = \sum_ {k = 0} ^ {n} \frac {1}{k !} \sup _ {0 \leqslant t \leqslant 1} | f ^ {(k)} (t) |,
\]

it is a Banach algebra (example 4 of I, p.~18). For every n, the space of characters \(\mathsf{X}(\mathsf{A}_{n})\) is identified with {[}0,1{]} and G under inclusion \(A_{n}\) into \(\mathcal{C}([0,1])\) (cf.~example 1 of I, p.~144).

\begin{enumerate}
\def\labelenumi{\arabic{enumi})}
\setcounter{enumi}{2}
\item
  Let \(\Delta\) be the disk of complex numbers \(z\) satisfying \(|z| \leqslant 1\) and let A be the complex Banach algebra of continuous functions on \(\Delta\) that are analytic in the interior of \(\Delta\) , equipped with the norm \(\|f\| = \sup_{z \in \Delta} |f(z)|\) (example 9 of I, p.~20). Then \(\mathsf{X}(\mathsf{A})\) is identified with \(\Delta\) and \(\mathcal{G}\) with the inclusion of A into \(\mathcal{C}(\Delta)\) (cf.~exerc. 6 of I, p.~193).
\item
  Consider the complex Banach algebra A of absolutely convergent Fourier series (example 8 of I, p.~19). For every element \(u\) of the unit circle U, the map \(f \mapsto f(u)\) is a character \(ev_{u}\) of A. If \(f_{0} \in A\) is the identity map of U, one has \(ev_{u}(f_{0}) = u\) , hence the map ev: \(u \mapsto ev_{u}\) from U into X(A) is injective; it is continuous. Let \(\chi \in \mathsf{X}(\mathsf{A})\) . We have \(\|f_{0}\| = \|f_{0}^{-1}\| = 1\) , hence \(|\chi(f_{0})| \leqslant 1\) and \(|\chi(f_{0})^{-1}| \leqslant 1\) . This shows that \(\chi(f_{0}) \in \mathbf{U}\) and there exists \(u \in U\) such that \(\chi(f_{0}) = \mathrm{ev}_{u}(f_{0})\) . Since \(\{f_{0}, f_{0}^{-1}\}\) topologically generates the unital algebra A, one has \(\chi = ev_{u}\) . Thus, the map ev is a homeomorphism from U onto X(A), by which one identifies these spaces. The Gelfand transformation of A is then identified with the inclusion of A into \(\mathcal{C}(\mathbf{U})\) .
\end{enumerate}

Since A is isomorphic to the Banach algebra \(\mathrm{L}^{1}(\mathbf{Z})\) (example 8 of I, p.~19), the space \(\mathsf{X}(\mathrm{L}^{1}(\mathbf{Z}))\) is identified with U and, for every element \((c_{n})\in\mathrm{L}^{1}(\mathbf{Z})\) , the Gelfand transform \(\mathcal{G}_{\mathrm{L}^{1}(\mathbf{Z})}((c_{n}))\) is identified with the function \(u\mapsto\sum_{n\in\mathbf{Z}}c_{n}u^{n}\) on U.

\begin{enumerate}
\def\labelenumi{\arabic{enumi})}
\setcounter{enumi}{4}
\tightlist
\item
  Let \(\Delta\) the unit disk of complex numbers \(z\) such that \(|z| \leqslant 1\) . Its boundary in \(\mathbf{C}\) is \(\mathbf{U}\) . Let A be the Banach algebra of complex functions \(f\) on \(\mathbf{U}\) such that there exists a continuous function \(\widetilde{f} \in \mathcal{C}(\Delta)\) extending \(f\) that is analytic in \(\mathring{\Delta}\) , equipped with the norm \(\|f\| = \sup_{z \in \mathbf{U}} |f(z)|\) . By virtue of the maximum principle (VAR, R1, p.~30, 3.3.7), one then has \(\|f\| = \sup_{z \in \Delta} |\widetilde{f}(z)|\) , and hence A coincides with the algebra of example 9 of I, p.~20. The set \(\mathsf{X}(\mathsf{A})\) is identified with \(\Delta\) and, if \(f \in \mathsf{A}\) , the map \(\mathcal{G}(f)\) is identified with the continuous extension of \(f\) into \(\Delta\) that is analytic in \(\mathring{\Delta}\) (cf.~exerc. 6 of I, p.~193).
\end{enumerate}

PROPOSITION 5. --- Let A be a commutative Banach algebra. For every \(x \in \mathbf{A}\) , the function \(\mathcal{G}(x)\) belongs to the commutative Banach algebra \(\mathcal{C}_0(\mathsf{X}(\mathbf{A}))\) of continuous functions on \(\mathsf{X}(\mathbf{A})\) vanishing at infinity.

By definition (cf.~no. 7 of I, p.~9), the function \(\mathcal{G}_{\mathrm{A}}^{\prime}(x)\colon \chi \mapsto \chi (x)\) is continuous on \(\mathsf{X}'(\mathsf{A})\) and vanishes at 0. Since \(\mathsf{X}'(\mathsf{A})\) is identified with the Alexandroff compactification of \(\mathsf{X}(\mathsf{A})\) according to cor. 1 of I, p.~29, the proposition follows.

PROPOSITION 6. --- Let A be a commutative Banach algebra and let \(x \in A\) .

\begin{enumerate}
\def\labelenumi{\alph{enumi})}
\item
  The union of the set of values of \(\mathcal{G}(x)\) and of \(\{0\}\) is equal to \(\mathrm{Sp}_{\mathrm{A}}^{\prime}(x)\) ;
\item
  If A has a unit element, the set of values of \(\mathcal{G}(x)\) is \(\mathrm{Sp}_{\mathrm{A}}(x)\) . In particular, for \(x\) to be invertible, it is necessary and sufficient that \(\mathcal{G}(x)\) does not vanish.
\end{enumerate}

Suppose that A has a unity element. We know that, for every \(\chi \in \mathsf{X}(\mathsf{A})\) , one has \(\chi(x) \in \mathrm{Sp}_{\mathsf{A}}(x)\) Conversely, let \(\lambda \in \mathrm{Sp}_{\mathsf{A}}(x)\) . Then \(x - \lambda\) is not invertible, hence belongs to a maximal ideal of A. There then exists \(\chi \in \mathsf{X}(\mathsf{A})\) such that \(\chi(x - \lambda) = 0\) (th. 2 of I, p.~30), whence \(b\) ).

Moving to the general case. Let \(\widetilde{\mathsf{A}}\) the Banach algebra obtained from A by adjoining an identity element; it is commutative. The set \(\mathrm{Sp}_{\widetilde{\mathsf{A}}}^{\prime}(x)\) is equal to \(\mathrm{Sp}_{\widetilde{\mathsf{A}}}^{\sim}(x)\) , that is, the set of values of \(\mathcal{G}_{\widetilde{\mathsf{A}}}^{\sim}(x)\) on \(\mathsf{X}(\widetilde{\mathsf{A}}) = \mathsf{X}'(\mathsf{A})\) Hence \(a\) ).

Example. --- Consider the Banach algebra A of absolutely convergent Fourier series (example 4). Prop. 6, b) implies that if \(\varphi\) is a function on the unit circle U admitting an absolutely convergent Fourier series, and if \(\varphi\) does not vanish, the function \(1/\varphi\) also admits an absolutely convergent Fourier series (“Wiener’s theorem”).

PROPOSITION 7. --- Let A be a commutative Banach algebra.

\begin{enumerate}
\def\labelenumi{\alph{enumi})}
\item
  The Gelfand transformation \(\mathcal{G}\) defines a morphism from A to \(\mathcal{C}_0(\mathsf{X}(\mathsf{A}))\) such that \(\| \mathcal{G}(x)\| = \varrho (x)\leqslant \| x\|\) for all \(x\in \mathrm{A}\) ;
\item
  For the Gelfand transform to \(\mathcal{G}\) it is isometric, it is necessary and sufficient that \(\|x^{2}\|=\|x\|^{2}\) for all \(x\in\mathrm{A}\) .
\end{enumerate}

The map \(\mathcal{G}\) is a morphism from A to \(\mathcal{C}_0(\mathsf{X}(\mathsf{A}))\) by no. 7 of I, p.~9 and prop. 5 of I, p.~37, and verifies \(\| \mathcal{G}(x)\| = \varrho (x)\) by prop. 6 and cor. 5 of I, p.~28. The assertion \(b)\) follows from \(a)\) and of Remark 1 of I, p.~21.

COROLLARY. --- Let A be a Banach algebra, and let x and y be commuting elements of A.

\begin{enumerate}
\def\labelenumi{\alph{enumi})}
\item
  We have \(\varrho(xy) \leqslant \varrho(x)\varrho(y)\) and \(\varrho(x + y) \leqslant \varrho(x) + \varrho(y)\) ;
\item
  If y is quasi-nilpotent, then \(\mathrm{Sp}_{\mathrm{A}}^{\prime}(x)=\mathrm{Sp}_{\mathrm{A}}^{\prime}(x+y)\) ; if moreover A is unital, then \(\mathrm{Sp}_{\mathrm{A}}(x)=\mathrm{Sp}_{\mathrm{A}}(x+y)\) .
\end{enumerate}

By considering the Banach algebra obtained from A by adjoining an identity element, one first reduces to the case where the algebra A is unital. Then, by considering the closed full subalgebra of A generated by x and y, one reduces to the case where A is commutative and unital. Assertion (a) is then a consequence of Prop. 7(a), and assertion (b) follows from Prop. 6 of I, p.~37 and Cor. 5 of I, p.~28.

PROPOSITION 8. --- Let A be a commutative Banach algebra. The following four sets are equal:

\begin{enumerate}
\def\labelenumi{(\roman{enumi})}
\item
  The kernel of the Gelfand transformation;
\item
  The set of elements x of A such that \(\mathrm{Sp}_{\mathrm{A}}^{\prime}(x)=\{0\}\) ;
\item
  The set of quasi-nilpotent elements of A;
\item
  The radical of A.
\end{enumerate}

We denote \(N_{1}\) , \(N_{2}\) , \(N_{3}\) , \(N_{4}\) , respectively, these sets. We have \(N_{1} = N_{2}\) (prop. 6, a)), and \(N_{2} = N_{3}\) (I, p.~28, cor. 5). By definition, the set \(N_{4}\) is the intersection of the regular maximal ideals of A; it is therefore the intersection of the kernels of the characters of A (th. 2 of I, p.~30), which is equal to \(N_{1}\) .

Remarks. --- 1) In general, the image of the Gelfand transformation is neither closed in \(\mathcal{C}_{0}(\mathsf{X}(\mathsf{A}))\) , nor dense in \(\mathcal{C}_{0}(\mathsf{X}(\mathsf{A}))\) (exerc. 7 of I, p.~193).

\begin{enumerate}
\def\labelenumi{\arabic{enumi})}
\setcounter{enumi}{1}
\item
  The image of the Gelfand transformation separates the points of \(\mathsf{X}(\mathsf{A})\) , since if \(\chi_1 \neq \chi_2\) are elements of \(\mathsf{X}(\mathsf{A})\) , there exists \(x \in \mathsf{A}\) such that \(\chi_1(x) \neq \chi_2(x)\) .
\item
  If \(\chi \in \mathsf{X}(\mathsf{A})\) , there exists an element of the image of the Gelfand transformation which does not vanish at \(\chi\) .
\item
  If A has an identity element, the image of the Gelfand transformation is a full subalgebra of the algebra of continuous functions on \(\mathsf{X}(\mathsf{A})\) (prop. 6, b)).
\end{enumerate}

Lemma 2. --- Let A be a commutative Banach algebra. Let M be a subset of X(A). Then M is closed for the Jacobson topology if and only if, for every character \(\chi \in \mathsf{X}(\mathsf{A}) = \mathsf{M}\) , there exists an element \(x\) of A such that \(\mathcal{G}(x)\) vanishes on M and is nonzero at \(\chi\) .

Let \(\Upsilon(\mathrm{M})\) the intersection of the kernels of the elements of M. The set M is closed for the Jacobson topology if and only if \(\mathrm{M} = \mathrm{V}(\Upsilon(\mathrm{M}))\) (cf.~I, p.~13 and I, p.~30). This condition is equivalent to saying that the elements \(\chi\) of M are precisely the characters which vanish on \(\Upsilon(\mathrm{M})\) . Consequently M is closed if and only if for every character \(\chi \notin \mathrm{M}\) , there exists \(x \in \Upsilon(\mathrm{M})\) such that \(\chi(x) \neq 0\) . This translates into \(\mathcal{G}(x)(\chi) \neq 0\) and \(\mathcal{G}(x)|_{\mathrm{M}} = 0\) .

\subsection{Morphisms of commutative Banach algebras}

PROPOSITION 9. --- Let A be a Banach algebra, and let B be a commutative semisimple Banach algebra. Every morphism from the underlying algebra of A into the underlying algebra of B is continuous.

Let \(h \colon \mathrm{A} \to \mathrm{B}\) an algebra homomorphism and let \((a, b) \in \mathrm{A} \times \mathrm{B}\) a point adherent to the graph \(\Gamma\) of \(h\) . Let \(\chi \in \mathsf{X}'(\mathrm{B})\) . The function \(x \mapsto \chi(h(x))\) of \(\mathrm{A}\) into \(\mathbf{C}\) is an algebra homomorphism, hence is continuous (I, p.~29, th. 1). The map from \(\mathrm{A} \times \mathrm{B}\) into \(\mathbf{C}\) given by \((x, y) \mapsto \chi(h(x)) - \chi(y)\) is then continuous, it vanishes on \(\Gamma\) therefore zero at \((a, b)\) We thus have \(\chi(h(a)) = \chi(b)\) for all \(\chi \in \mathsf{X}'(\mathrm{B})\) . Since B is semisimple, we have \(h(a) = b\) . Thus the graph of \(h\) is closed and therefore \(h\) is continuous (EVT, I, p.~19, cor. 5).

COROLLARY. --- On a commutative complex algebra without radical, two norms defining Banach algebra structures are equivalent.

It suffices to apply Prop. 9 to the identity map of the algebra.

Let A and B be commutative Banach algebras. By No. 7 of I, p.~9, if \(h\colon A\to B\) is a surjective morphism, \(\mathsf{X}'(h)\) is a homeomorphism from \(\mathsf{X}'(\mathsf{B})\) on a closed subspace of \(\mathsf{X}'(\mathsf{A})\) which maps 0 to 0. (In the case where \(h\) is the injection of a subalgebra into \(\mathcal{C}(\mathsf{X})\) , the map \(\mathsf{X}'(h)\) is injective under much weaker hypotheses, cf.~I, p.~142, prop. 1, d).)

Now let \(h \colon \mathrm{A} \to \mathrm{B}\) an injective morphism. In general, \(\mathsf{X}'(h)\) is not surjective, but the following proposition provides a necessary condition for this to be the case.

PROPOSITION 10. --- Let A and B be unital commutative Banach algebras, \(h\colon A\to B\) a unital algebra homomorphism, not necessarily continuous. If \(\mathsf{X}(h)\) is surjective, then \(h(\mathrm{A})\) is a full subalgebra of B.

Let \(x \in A\) such that \(h(x)\) be invertible in B. For every \(\chi \in \mathsf{X}(\mathsf{A})\) , there exists \(\xi \in \mathsf{X}(\mathsf{B})\) such that \(\chi = \mathsf{X}(h)(\xi)\) , hence \(\chi(x) = \xi(h(x)) \neq 0\) . Prop. 6 of I, p.~37 then shows that x is invertible in A, and therefore that \(h(x)\) is invertible in \(h(\mathsf{A})\) .

The necessary condition of the proposition is not sufficient, even if h is isometric (I, p.~168, exerc. 14). We nevertheless have the following result:

PROPOSITION 11. --- Let A and B be commutative unital Banach algebras, let a be an element of A and let h: A → B be an injective unital morphism (not necessarily continuous). Suppose that the closed full subalgebra of A generated by a is equal to A. The following conditions are equivalent:

\begin{enumerate}
\def\labelenumi{(\roman{enumi})}
\item
  \(\mathsf{X}(h)\) is surjective;
\item
  \(h(\mathrm{A})\) is a full subalgebra of B;
\item
  \(\mathrm{Sp}_{\mathrm{A}}(a)=\mathrm{Sp}_{\mathrm{B}}(h(a)).\)
\item
  \(\Longrightarrow\) (ii) follows from prop. 10.
\item
  \(\Longrightarrow\) (iii) follows from the formula \(\mathrm{Sp}_{\mathrm{A}}(a) = \mathrm{Sp}_{h(\mathrm{A})}(h(a)) = \mathrm{Sp}_{\mathrm{B}}(h(a))\) , valid since \(h(\mathrm{A})\) is a full subalgebra of B.
\item
  \(\Longrightarrow\) (i) by formula (4) of no. 6 of I, p.~6, we have the commutative diagram
\end{enumerate}

\[
\begin{array}{c} \mathsf {X} (\mathrm{B}) \xrightarrow {\mathsf {X} (h)} \mathsf {X} (\mathrm{A}) \\ \Big \downarrow^ {\mathcal {G} _ {\mathrm{B}} (h (a))} \qquad \qquad \qquad \Big \downarrow^ {\mathcal {G} _ {\mathrm{A}} (a)} \\ \operatorname{Sp} _ {\mathrm{B}} (h (a)) \xrightarrow {i} \operatorname{Sp} _ {\mathrm{A}} (a) \end{array}
\]

where the vertical arrows denote surjective maps (I, p.~37, prop. 6) and \(i\) is the canonical inclusion. Hypothesis (iii) means that \(i\) is bijective. Moreover, the surjective map \(\mathcal{G}_{\mathrm{A}}(a)\colon \mathsf{X}(\mathrm{A})\to \mathrm{Sp}_{\mathrm{A}}(a)\) is bijective: indeed, for every characters \(\chi_{1}\) and \(\chi_{2}\) of A, the set \(\mathrm{A}_{\chi_1,\chi_2} = \{x\in \mathrm{A}\mid \chi_1(x) = \chi_2(x)\}\) is a closed full subalgebra of A. By hypothesis, we therefore have \(\mathrm{A}_{\chi_1,\chi_2} = \mathrm{A}\) if \(a\in \mathrm{A}_{\chi_1,\chi_2}\) , that is, \(\chi_{1} = \chi_{2}\) if \(\mathcal{G}_{\mathrm{A}}(a)\chi_{1} = \mathcal{G}_{\mathrm{A}}(a)\chi_{2}\) . The diagram then implies that the map \(\mathsf{X}(h)\) is surjective.

\subsection{Simultaneous spectrum}

Let \(\Lambda\) a set. Let \(C_{\Lambda} = \mathbf{C}[(X_{\lambda})_{\lambda \in \Lambda}]\) be the unital algebra of complex polynomials in a family of indeterminates \((X_{\lambda})_{\lambda \in \Lambda}\) . For every \(\chi \in X(C_{\Lambda})\) , one has \((\chi (X_{\lambda}))_{\lambda \in \Lambda} \in C^{\Lambda}\) ; the map \(\chi \mapsto (\chi (X_{\lambda}))_{\lambda \in \Lambda}\) is a homeomorphism from \(X(C_{\Lambda})\) onto the product space \(C^{\Lambda}\) , by which these spaces are identified.

On the other hand, let A be a commutative unital Banach algebra and let \(x = (x_{\lambda})_{\lambda \in \Lambda}\) be a family of elements of A. There exists a unique unital morphism h from \(C_{\Lambda}\) to A such that \(h(\mathrm{X}_{\lambda}) = x_{\lambda}\) for all \(\lambda\) . The continuous map \(\mathsf{X}(h)\) of \(\mathsf{X}(\mathrm{A})\) into \(C^{\Lambda}\) is the map which to \(\chi\) associates the family \((\chi(x_{\lambda}))_{\lambda \in \Lambda}\) . It is called the map of \(\mathsf{X}(\mathrm{A})\) into \(C^{\Lambda}\) defined by x.

DEFINITION 2. --- The image of the map X(h) is called the simultaneous spectrum of x, and is denoted \(\mathrm{Sp}_{\mathrm{A}}^{\Lambda}(x)\) or \(\mathrm{Sp}^{\Lambda}(x)\) .

The simultaneous spectrum of \(x\) is a compact subset of \(\mathbf{C}^{\Lambda}\) . An element \(c = (c_{\lambda}) \in \mathbf{C}^{\Lambda}\) belongs to \(\mathrm{Sp}_{\mathrm{A}}^{\Lambda}(x)\) if and only if the elements \(x_{\lambda} - c_{\lambda}\) belong to a same maximal ideal of \(\mathrm{A}\) , in other words if the family \((x_{\lambda} - c_{\lambda})_{\lambda \in \Lambda}\) does not generate the algebra \(\mathrm{A}\) .

If \(\Lambda\) contains a single element, so that the family \(x\) reduces to a single element \(x \in \mathsf{A}\) , one has \(\mathrm{Sp}_{\mathsf{A}}^{\Lambda}(x) = \mathrm{Sp}_{\mathsf{A}}(x)\) (I, p.~37, prop. 6, \(b\) )). If \(\Lambda' \subset \Lambda\) , then \(\mathrm{Sp}_{\mathsf{A}}^{\Lambda'}((x_{\lambda})_{\lambda \in \Lambda'})\) is the image of \(\mathrm{Sp}_{\mathsf{A}}^{\Lambda}((x_{\lambda})_{\lambda \in \Lambda})\) by the canonical projection map of \(\mathbf{C}^{\Lambda}\) on \(\mathbf{C}^{\Lambda'}\) . In particular, we have

\[
\operatorname{Sp} _ {\mathrm{A}} ^ {\Lambda} (x) \subset \prod_ {\lambda \in \Lambda} \operatorname{Sp} _ {\mathrm{A}} (x _ {\lambda}).
\]

We denote \(z_{\lambda}\) , for \(\lambda \in \Lambda\) , the coordinate functions on \(\mathbf{C}^{\Lambda}\) . If \(\chi \in \mathsf{X}(\mathsf{A})\) , the value at \(\chi\) of \(z_{\lambda} \circ \mathsf{X}(h)\) is \(\chi(x_{\lambda})\) , hence \(z_{\lambda} \circ \mathsf{X}(h) = \mathcal{G}(x_{\lambda})\) .

Let A and B be commutative unital Banach algebras, \(\varphi\) a unital morphism from A to B, and \(x = (x_{\lambda})_{\lambda \in \Lambda}\) a family of elements of A. Let \(\varphi(x)\) the family \((\varphi(x_{\lambda}))_{\lambda \in \Lambda}\) of elements of B. We have, for every \(\chi \in \mathsf{X}(\mathsf{B})\) , and every \(\lambda \in \Lambda\)

\[
\chi (\varphi (x _ {\lambda})) = (\mathsf {X} (\varphi) (\chi)) (x _ {\lambda}),
\]

hence \(\operatorname{Sp}_{\mathrm{B}}^{\Lambda}(\varphi(x)) \subset \operatorname{Sp}_{\mathrm{A}}^{\Lambda}(x)\) . The diagram

\[
\begin{array}{c} \mathsf {X} (\mathrm{B}) \xrightarrow {\mathsf {X} (\varphi)} \mathsf {X} (\mathrm{A}) \\ \Big \downarrow \qquad \qquad \qquad \qquad \Big \downarrow \\ \operatorname{Sp} _ {\mathrm{B}} ^ {\Lambda} (\varphi (x)) \xrightarrow {i} \operatorname{Sp} _ {\mathrm{A}} ^ {\Lambda} (x) \end{array}\tag{1}
\]

where i denotes the inclusion, and the vertical arrows denote the maps defined by the families \(\varphi(x)\) and x, is therefore commutative.

Example. --- Let \(K \subset C^{\Lambda}\) a compact subset. Let \(z = (z_{\lambda})_{\lambda \in \Lambda}\) the family in \(\mathcal{C}(K)\) of the restrictions to K of the coordinate functions of \(\mathbf{C}^{\Lambda}\) . Then the simultaneous spectrum \(\mathrm{Sp}_{\mathcal{C}(\mathrm{K})}^{\Lambda}(z)\) is equal to K. Indeed, by cor. 2 of I, p.~32, every character \(\chi\) of \(\mathcal{C}(\mathrm{K})\) is of the form \(f \mapsto f(x)\) for an element \(x \in \mathrm{K}\) , and one then has \((\chi(z_{\lambda}))_{\lambda \in \Lambda} = x\) .

PROPOSITION 12. — Let \(\Lambda\) a set and A a commutative unital Banach algebra. Let \(x = (x_{\lambda})_{\lambda \in \Lambda}\) be a family of elements of A.

\begin{enumerate}
\def\labelenumi{\alph{enumi})}
\item
  We assume that the full subalgebra of A generated by the family x is dense in A. The map from X(A) to \(\mathbf{C}^{\Lambda}\) defined by x is a homeomorphism from X(A) onto the simultaneous spectrum \(\mathrm{Sp}_{\mathrm{A}}^{\Lambda}(x)\) ;
\item
  We assume that the unital subalgebra of A generated by the family x is dense in A. For every \(c \in C^{\Lambda}\) , the following conditions are equivalent:
\end{enumerate}

\begin{enumerate}
\def\labelenumi{(\roman{enumi})}
\item
  \(c \in \mathrm{Sp}_{\mathrm{A}}^{\Lambda}(x)\) ;
\item
  \(|\mathrm{P}(c)| \leqslant \varrho(\mathrm{P}(x))\) for every polynomial \(\mathrm{P} \in \mathbf{C}[(\mathrm{X}_{\lambda})_{\lambda \in \Lambda}]\) ;
\item
  \(|\mathrm{P}(c)| \leqslant \| \mathrm{P}(x)\|\) for every polynomial \(\mathrm{P} \in \mathbf{C}[(\mathrm{X}_{\lambda})_{\lambda \in \Lambda}]\) .
\end{enumerate}

\begin{enumerate}
\def\labelenumi{\alph{enumi})}
\item
  The map from \(\mathsf{X}(\mathsf{A})\) into \(\mathrm{Sp}_{\mathrm{A}}^{\Lambda}(x)\) defined by the family \(x\) is continuous and surjective. Let \(\chi\) , \(\chi' \in \mathsf{X}(\mathsf{A})\) be characters having the same image, that is, such that \(\chi(x_{\lambda}) = \chi'(x_{\lambda})\) for all \(\lambda \in \Lambda\) . The characters \(\chi\) and \(\chi'\) coincide on elements of the form \(\mathrm{P}(x)\mathrm{Q}(x)^{-1}\) , where \(\mathrm{P} \in \mathbf{C}[(\mathrm{X}_{\lambda})]\) , \(\mathrm{Q} \in \mathbf{C}[(\mathrm{X}_{\lambda})]\) and \(\mathrm{Q}(x)\) is invertible in A, that is, on the full subalgebra of A generated by the elements \(x_{\lambda}\) (lemma 2 of I, p.~6). Since \(\chi\) and \(\chi'\) are continuous (th. 1 of I, p.~29), they are therefore equal on A. This proves that \(\mathsf{X}(h)\) is a continuous bijection from \(\mathsf{X}(\mathsf{A})\) on \(\mathrm{Sp}_{\mathrm{A}}^{\Lambda}(x)\) , and consequently it is a homeomorphism since \(\mathsf{X}(\mathsf{A})\) is compact.
\item
  Let us show that (i) implies (ii): if \(c = (c_{\lambda})_{\lambda \in \Lambda} \in \mathrm{Sp}_{\mathrm{A}}^{\Lambda}(x)\) , there exists \(\chi \in \mathsf{X}(\mathsf{A})\) such that \(c_{\lambda} = \chi(x_{\lambda})\) for all \(\lambda\) . For every \(\mathrm{P} \in \mathbf{C}[(\mathrm{X}_{\lambda})]\) , we therefore have \(|\mathrm{P}(c)| = |\mathrm{P}((\chi(x_{\lambda}))_{\lambda \in \Lambda})| = |\chi(\mathrm{P}(x))| \leqslant \varrho(\mathrm{P}(x))\) .
\end{enumerate}

Assertion (ii) implies (iii) by virtue of the inequality \(\varrho(x) \leqslant \|x\|\) , valid for all \(x \in A\) by definition of the spectral radius.

Let us finally show that (iii) implies (i). Let \(c = (c_{\lambda})_{\lambda \in \Lambda} \in \mathbf{C}^{\Lambda}\) such that

\[
| \mathrm{P} (c) | \leqslant \| \mathrm{P} (x) \|\tag{2}
\]

for all \(P \in \mathbf{C}[(X_{\lambda})]\) . Let A' be the unital subalgebra of A generated by the family x; its elements are of the form \(P(x)\) for \(P \in \mathbf{C}[(X_{\lambda})]\) . The estimate (2) implies that the condition \(P(x) = 0\) implies \(\mathrm{P}(c)=0\) There therefore exists a unital algebra morphism \(\xi\) of \(A'\) into C such that \(\xi(x_{\lambda})=c_{\lambda}\) for all \(\lambda\in\Lambda\) . By (2), the morphism \(\xi\) is continuous. It therefore extends by continuity to a character \(\chi\) of \(\overline{\mathrm{A}'}=\mathrm{A}\) , which satisfies \(c=(\chi(x_{\lambda}))_{\lambda\in\Lambda}\in\mathrm{Sp}_{\mathrm{A}}^{\Lambda}(x)\) . This completes the proof.

\subsection{Polynomially convex sets}

DEFINITION 3. --- Let \(\Lambda\) a set and V a subset of \(\mathbf{C}^{\Lambda}\) . We say that V is polynomially convex if V is the set of points \((c_{\lambda})_{\lambda \in \Lambda}\) of \(\mathbf{C}^{\Lambda}\) such that

\[
| \mathrm{P} ((c _ {\lambda})) | \leqslant \sup _ {c \in \mathrm{V}} | \mathrm{P} (c) |
\]

for all \(\mathrm{P}\in \mathbf{C}[(\mathrm{X}_{\lambda})]\)

Lemma 3. --- Let \(\Lambda\) a set. A subset V of \(\mathbf{C}^{\Lambda}\) is polynomially convex if and only if there exist a family \((\mathrm{P}_i)_{i\in \mathrm{I}}\) of elements of \(\mathbf{C}[(\mathrm{X}_{\lambda})_{\lambda \in \Lambda}]\) and a family \((\mathrm{M}_i)_{i\in \mathrm{I}}\) of elements of \([0, +\infty]\) such that V is the set of \(c\in \mathbf{C}^{\Lambda}\) satisfying

\[
| \mathrm{P} _ {i} (c) | \leqslant \mathrm{M} _ {i}
\]

for all i ∈ I.

If the subset V of \(C^{\Lambda}\) is polynomially convex, it satisfies the above condition for the family consisting of the elements P of \(\mathbf{C}[(X_{\lambda})_{\lambda\in\Lambda}]\) by setting \(M_{P}=\sup_{c\in V}|P(c)|\) .

Conversely, let \((\mathrm{P}_{i})_{i\in\mathrm{I}}\) be a family of elements of \(\mathbf{C}[(\mathrm{X}_{\lambda})_{\lambda\in\Lambda}]\) and \((\mathrm{M}_{i})_{i\in\mathrm{I}}\) be a family of elements of \([0,+\infty]\) . Let V be the set of c in \(C^{\Lambda}\) such that \(|\mathrm{P}_{i}(c)|\leqslant\mathrm{M}_{i}\) for all \(i\in I\) . We then have \(\sup_{c\in\mathrm{V}}|\mathrm{P}_{i}(c)|\leqslant\mathrm{M}_{i}\) for \(i\in I\) . Suppose that \(x\in C^{\Lambda}\) satisfies

\[
| \mathrm{P} (x) | \leqslant \sup _ {c \in \mathrm{V}} | \mathrm{P} (c) |
\]

for all \(\mathrm{P}\in\mathbf{C}[(\mathrm{X}_{\lambda})]\) . For \(i\in I\) , we have in particular \(|\mathrm{P}_{i}(x)|\leqslant\mathrm{M}_{i}\) , hence \(x\in V\) . Conversely, for every element \(x\in V\) and every polynomial \(\mathrm{P}\in\mathbf{C}[(\mathrm{X}_{\lambda})]\) , one has \(|\mathrm{P}(x)|\leqslant\sup_{c\in\mathrm{V}}|\mathrm{P}(c)|\) Consequently, the set V is polynomially convex.

Lemma 4. --- Let A be a unital commutative Banach algebra. Let \(\Lambda\) a set and \(x = (x_{\lambda})_{\lambda \in \Lambda}\) a family of elements of A. If the unital subalgebra generated by the family x is dense in A, then the simultaneous spectrum \(\mathrm{Sp}_{\mathrm{A}}^{\Lambda}(x)\) is polynomially convex.

This follows from the assertion. \(b\) ) of prop. 12 of I, p.~43 and of definition 3.

Any intersection of polynomially convex subsets of \(C^{\Lambda}\) is polynomially convex (Lemma 3). This justifies the following definition:

DEFINITION 4. --- Let \(\Lambda\) a set and V a subset of \(\mathbf{C}^{\Lambda}\) The polynomially convex hull of V is the smallest polynomially convex subset of \(\mathbf{C}^{\Lambda}\) containing V.

The polynomially convex hull of V is the set of c belonging to \(C^{\Lambda}\) such that \(|P(c)| \leqslant \sup_{W} |P|\) for all \(P \in \mathbf{C}[(X_{\lambda})]\) Indeed, this set is polynomially convex by Lemma 3, and is contained in every polynomially convex set containing V by definition.

Example. --- Let \(\Lambda\) a finite set and \(V \subset \mathbf{C}^{\Lambda}\) a compact convex subset. Then \(V\) is polynomially convex. Indeed, let \(W\) be the polynomially convex hull of \(V\) . Let us prove that \(W \subset V\) , which will imply the assertion. Let \(x \in \mathbf{C}^{\Lambda} - V\) . There exists a real hyperplane \(H\) into \(\mathbf{C}^{\Lambda}\) which strictly separates \(x\) and \(V\) (EVT, II, p.~41, prop. 4). Let \(f_{\mathbf{R}}\) a form \(\mathbf{R}\) -linear on \(\mathbf{C}^{\Lambda}\) and \(\alpha \in \mathbf{R}\) such that \(H\) is the set of \(y \in \mathbf{C}^{\Lambda}\) satisfying \(f_{\mathbf{R}}(y) = \alpha\) . Let \(f\) a linear form on \(\mathbf{C}^{\Lambda}\) such that \(f_{\mathbf{R}} = \mathcal{R}(f)\) . We therefore have

\[
\mathcal {R} (f (x)) > \sup _ {y \in V} \mathcal {R} (f (y)).
\]

For every \(t \in R\) and \(y \in V\) , set \(f_{t}(y) = t + f(y)\) . We have \(|f_{t}| - \mathcal{R}(f_{t}) \to 0\) into \(\mathcal{C}(\mathbf{C}^{\Lambda}, \mathbf{R})\) endowed with the topology of compact convergence when \(t \to +\infty\) . For t sufficiently large, it follows that \(|f_{t}(x)| > \sup_{y \in V} |f_{t}(y)|\) since V is compact. Thus \(x \in C^{\Lambda} - W\) since \(f_{t}\) is a polynomial function.

Lemma 5. --- Let K be a compact subset of C. We denote by \(\widehat{\mathbf{K}}\) the union of K and the connected components of \(\mathbf{C} - \mathbf{K}\) which are relatively compact. Then the set \(\widehat{\mathbf{K}}\) is compact.

Since K is compact, there exists a real number r > 0 such that K is contained in the open disk D centered at 0 with radius r. Then C-K contains C-D. Since the space C-D is connected (being homeomorphic to \([r, +\infty[\times\mathbf{S}^1)\) , it is contained in a connected component U of \(\mathbf{C}-\mathbf{K}\) . Every other connected component of \(\mathbf{C}-\mathbf{K}\) is contained in D, hence is bounded. The connected component U is then the unique unbounded connected component of \(\mathbf{C}-\mathbf{K}\) , that is, \(\mathrm{U} = \mathbf{C}-\widehat{\mathrm{K}}\) . Since U is open and contains the complement of the disk D, the set \(\widehat{\mathrm{K}}\) is compact.

PROPOSITION 13. --- Let \(n \geqslant 1\) be an integer. Let \(\mathbf{K} \subset \mathbf{C}^{n}\) be a closed subset and V its polynomially convex hull.

\begin{enumerate}
\def\labelenumi{\alph{enumi})}
\item
  Every bounded connected component of \(C^{n} - K\) is contained in V;
\item
  If n = 1, and if K is compact, then V is the union \(\widehat{K}\) of K and the bounded connected components of C-K.
\end{enumerate}

Since \(K \subset C^{n}\) is closed, its complement \(C^{n}-K\) is open, hence locally connected, so that each connected component of \(C^{n}-K\) is open. The maximum principle (VAR, R1, p.~29, 3.3.7) then shows that every bounded connected component of \(C^{n}-K\) is contained in V, which proves assertion a).

Now suppose \(n = 1\) and K compact. Let \(\widehat{\mathbf{K}}\) denote the union of K and the bounded connected components of \(\mathbf{C} - \mathbf{K}\) , so that \(\widehat{\mathbf{K}} \subset \mathbf{V}\) ; by the preceding, the set \(\widehat{\mathbf{K}}\) is compact (lemma 5). Let A be the commutative unital Banach algebra \(\mathcal{C}(\widehat{\mathbf{K}})\) . Let \(x \in \mathbf{A}\) the identity function on \(\widehat{\mathbf{K}}\) and B the closed unital subalgebra of A generated by \(x\) . We have \(\mathrm{Sp}_{\mathrm{A}}(x) = \widehat{\mathbf{K}}\) (example 3 of I, p.~17), hence \(\mathbf{C} - \mathrm{Sp}_{\mathrm{A}}(x)\) is connected. Consequently we have \(\mathrm{Sp}_{\mathrm{B}}(x) = \mathrm{Sp}_{\mathrm{A}}(x)\) (cor. of prop. 6 of I, p.~28). Since \(\mathrm{Sp}_{\mathrm{B}}(x)\) is polynomially convex by lemma 4, this shows that \(\widehat{\mathbf{K}}\) is polynomially convex and therefore that \(\mathbf{V} \subset \widehat{\mathbf{K}}\) .

The second part of the proposition does not extend to the case \(n \geqslant 2\) (cf.~exerc. 23 of I, p.~170).

PROPOSITION 14. --- Let \(\Lambda\) be a set. Let A be a commutative unital Banach algebra, \(x = (x_{\lambda})_{\lambda \in \Lambda}\) a family of elements of A and A' the unital Banach subalgebra generated by \(x\) . Then the simultaneous spectrum \(\mathrm{Sp}_{\mathrm{A}'}^{\Lambda}(x)\) is the polynomially convex hull of \(\mathrm{Sp}_{\mathrm{A}}^{\Lambda}(x)\) .

Indeed, prop. 12 of I, p.~43, b) proves that \(\mathrm{Sp}_{\mathrm{A}'}^{\Lambda}(x)\) is the set of \(c \in \mathbf{C}^{\Lambda}\) such that \(|\mathrm{P}(c)| \leqslant \varrho(\mathrm{P}(x))\) for all \(\mathrm{P} \in \mathbf{C}[(\mathrm{X}_{\lambda})]\) .

Now we have

\[
\varrho (\mathrm{P} (x)) = \sup _ {\chi \in \mathsf {X} (\mathrm{A})} | \chi (\mathrm{P} (x)) | = \sup _ {\chi \in \mathsf {X} (\mathrm{A})} | \mathrm{P} ((\chi (x _ {\lambda})) _ {\lambda \in \Lambda}) | = \sup _ {c \in \operatorname{Sp} _ {\mathrm{A}} ^ {\Lambda} (x)} | \mathrm{P} (c) |,
\]

by prop. 7 of I, p.~38, a) and prop. 12 of I, p.~43, a). The result then follows from lemma 3.

PROPOSITION 15. --- Let \(\Lambda\) be a set. Let K be a compact polynomially convex subset of \(\mathbf{C}^{\Lambda}\) . Let \(\mathrm{A}_0\) the set of restrictions to K of polynomial functions on \(\mathbf{C}^{\Lambda}\) , and let A be the closure of \(\mathrm{A}_0\) in algebra \(\mathcal{C}(\mathrm{K})\) .

Let ev: K → X(A) be the map defined by x → ev\_x, where ev\_x is the character f → f(x) of A. Let z = (z\_λ)\_\{λ∈Λ\} be the family in A of restrictions to K of the coordinate functions on C\^{}Λ, and let φ be the surjective map from X(A) to Sp\_A\^{}Λ(z) defined by the family z.

Then we have K = Sp \(_{A}^{\Lambda}(z)\) , and the maps ev: K → X(A) and φ: X(A) → K are reciprocal homeomorphisms.

The map \(\varphi \circ\) ev is the identity map of K. In particular, K is contained in the image \(\mathrm{Sp}_{\mathrm{A}}^{\Lambda}(z)\) of \(\varphi\) .

Prop. 14 implies that \(\mathrm{Sp}_{\mathrm{A}}^{\Lambda}(z)\) is the polynomially convex hull of \(\mathrm{Sp}_{\mathcal{C}(\mathrm{K})}^{\Lambda}(z)\) . Since, on the other hand, we have \(\mathrm{Sp}_{\mathcal{C}(\mathrm{K})}^{\Lambda}(z) = \mathrm{K}\) (I, p.~42, example), which is polynomially convex by hypothesis, we deduce that \(\mathrm{Sp}_{\mathrm{A}}^{\Lambda}(z) = \mathrm{K}\) .

Since the unital subalgebra generated by the family of elements \(z_{\lambda}\) is dense in A, the map \(\varphi\) is a homeomorphism from \(\mathsf{X}(\mathsf{A})\) on \(\mathrm{Sp}_{\mathsf{A}}^{\Lambda}(z)=\mathrm{K}\) Proposition 12 of Book I, p. 43(a). The identity \(\varphi\circ ev=Id_{K}\) then shows that ev is the inverse homeomorphism of \(\varphi\) .

For every subset \(\Lambda'\) of \(\Lambda\) , we denote by \(\mathrm{pr}_{\Lambda'}\) the canonical projection \(\mathbf{C}^{\Lambda} \to \mathbf{C}^{\Lambda'}\) . Let W be a subset of \(\mathbf{C}^{\Lambda}\) and let V be its polynomially convex hull. Set \(W' = \mathrm{pr}_{\Lambda'} W\) . Since every element of \(\mathbf{C}[(X_{\lambda})_{\lambda \in \Lambda'}]\) is identified with an element of \(\mathbf{C}[(X_{\lambda})_{\lambda \in \Lambda}]\) the polynomially convex hull of \(W'\) is contained in \(\mathrm{pr}_{\Lambda'} V\) .

Lemma 6. --- Let \(K \subset C^{\Lambda}\) a compact polynomially convex subset and U a neighborhood of K. There exists a finite subset \(\Lambda_{0}\) of \(\Lambda\) such that, for every subset \(\Lambda'\) of \(\Lambda\) containing \(\Lambda_{0}\) , the set \(\mathrm{pr}_{\Lambda'}(\mathrm{U})\) contains the polynomially convex hull of \(\mathrm{pr}_{\Lambda'}(\mathrm{K})\) .

Since K is compact, there exists a family of compact disks \(D_{\lambda}\) in C with center 0 and radii \(R_{\lambda}\) such that K is contained in the product \(\mathrm{D} = \prod_{\lambda} \mathrm{D}_{\lambda}\) . For every \(\mathrm{P} \in \mathbf{C}[(\mathrm{X}_{\lambda})]\) , let \(\mathrm{K}_{\mathrm{P}}\) the set of \(x \in \mathbf{C}^{\Lambda}\) such that

\[
| \mathrm{P} (x) | \leqslant \sup _ {c \in \mathrm{K}} | \mathrm{P} (c) |.
\]

We have

\[
\mathrm{D} \cap \bigcap_ {\mathrm{P}} \mathrm{K} _ {\mathrm{P}} = \mathrm{K}\tag{3}
\]

which is contained in U. Thus, the space D is the union of the open set \(D \cap U\) and of the family of open sets \(D - K_{P}\) for \(P \in \mathbf{C}[(X_{\lambda})]\) Since D is compact, there exists an integer \(q \geqslant 1\) and polynomials \(P_{1}, \ldots, P_{q} \in \mathbf{C}[(X_{\lambda})]\) such that:

\[
\mathrm{D} \cap \mathrm{K} _ {\mathrm{P} _ {1}} \cap \dots \cap \mathrm{K} _ {\mathrm{P} _ {q}} \subset \mathrm{U}.\tag{4}
\]

There exists a finite set \(\Lambda_0 \subset \Lambda\) such that \(P_i \in \mathbf{C}[(X_\lambda)_{\lambda \in \Lambda_0}]\) for \(1 \leqslant i \leqslant q\) . Let us show that \(\Lambda_0\) has the required property.

Let \(\Lambda'\) a subset of \(\Lambda\) containing \(\Lambda_0\) . Let E be the subset of \(\mathbf{C}^{\Lambda'}\) consisting of the elements \(c = (c_{\lambda})_{\lambda \in \Lambda'}\) defined by the inequalities \(|c_{\lambda}| \leqslant \mathrm{R}_{\lambda}\) for \(\lambda \in \Lambda'\) , and

\[
| \mathrm{P} _ {i} (c) | \leqslant \sup _ {c \in \mathrm{K}} | \mathrm{P} _ {i} (c) |
\]

for \(i = 1, \ldots, q\) . The set E is polynomially convex (lemma 3), and formula (3) shows that \(\mathrm{pr}_{\Lambda'}(\mathrm{K}) \subset \mathrm{E}\) .

On the other hand, let \(c = (c_{\lambda})_{\lambda \in \Lambda'} \in \mathrm{E}\) ; let \(d = (d_{\lambda})_{\lambda \in \Lambda}\) be the element of \(\mathbf{C}^{\Lambda}\) defined by \(d_{\lambda} = c_{\lambda}\) for \(\lambda \in \Lambda'\) and \(d_{\lambda} = 0\) for \(\lambda \in \Lambda - \Lambda'\) Then (4) implies that \(d \in \mathrm{U}\) and therefore that \(c \in \mathrm{pr}_{\Lambda'}(\mathrm{U})\) Thus, \(\mathrm{E} \subset \mathrm{pr}_{\Lambda'}(\mathrm{U})\) , which completes the proof.

Lemma 7. --- Let \(n \geqslant 1\) an integer and K a polynomially convex compact subset of \(\mathbf{C}^{n}\) Then K admits a fundamental system of compact neighborhoods that are polynomially convex.

There exists a compact polydisk (cf.~VAR, R1, p.~24) \(\Delta\) of \(\mathbf{C}^n\) which is a neighborhood of K. Since K is polynomially convex, there exists a family \((\mathrm{P}_i)_{i\in \mathrm{I}}\) of elements of \(\mathbf{C}[\mathrm{X}_1,\dots ,\mathrm{X}_n]\) , and a family \((\mathrm{M}_i)_{i\in \mathrm{I}}\) of positive real numbers such that K is the set of \(z\in\) \(\Delta\) satisfying \(|\mathrm{P}_i(z)|\leqslant \mathrm{M}_i\) for all \(i\) (Lemma 3). For every finite subset J of I and every \(\varepsilon >0\) , let \(\mathrm{K}_{\mathrm{J},\varepsilon}\) the set of \(z\in \Delta\) such that \(|\mathrm{P}_i(z)|\leqslant \mathrm{M}_i + \varepsilon\) for \(i\in \mathrm{J}\) . Then each set \(\mathrm{K}_{\mathrm{J},\varepsilon}\) is a compact polynomially convex neighborhood of K (loc. cit.), and the intersection of the sets \(\mathrm{K}_{\mathrm{J},\varepsilon}\) is K. The sets \(\mathrm{K}_{\mathrm{J},\varepsilon}\) therefore form a fundamental system of polynomially convex neighborhoods of K (TG, I, p.~60, th. 1).

\section{HOLOMORPHIC FUNCTIONAL CALCULUS}

\subsection{Germs of holomorphic functions}

Let E and F be complex Banach spaces. We recall (cf. VAR, R1, p.~26, 3.2.1, p.~22, 3.1 and p.~88, App.) that a holomorphic mapping defined on an open subset U of E and taking values in F is a mapping \(f \colon \mathrm{U} \to \mathrm{F}\) such that, for all \(x \in \mathrm{U}\) there exists a convergent series

\[
f _ {x} = \sum_ {k \geqslant 0} f _ {x, k}
\]

satisfying \(f(x + y) = f_x(y)\) for all \(y \in E\) sufficiently close to 0, where \(f_{x,k} \colon E \to C\) is a continuous homogeneous polynomial of degree k on E with values in F, that is, a mapping of the form

\[
f _ {x, k} (y) = \widetilde {f} _ {x, k} (y, \ldots , y)
\]

where \(\widetilde{f}_{x,k}\colon\mathrm{E}^{k}\to\mathrm{F}\) is a map \(k\) continuous multilinear. We denote \(\mathcal{O}(\mathrm{U};\mathrm{F})\) the complex vector space of holomorphic functions on U with values in F, equipped with the topology of compact convergence. It is a locally convex topological vector space, whose topology is defined by the seminorms \(f\mapsto\sup_{z\in\mathrm{K}}\|f(z)\|\) where K ranges over the set of compact subsets of U.

Let G be a complex Banach space and V an open subset of G. For every holomorphic map \(\varphi\colon\mathrm{V}\to\mathrm{U}\) , the map \(\varphi^{*}\colon f\mapsto f\circ\varphi\) is a continuous linear map from \(\mathcal{O}(\mathrm{U};\mathrm{F})\) into \(\mathcal{O}(\mathrm{V};\mathrm{F})\) .

If H is a complex Banach space and \(\varphi\colon\mathrm{F}\to\mathrm{H}\) a continuous linear map, then the map \(f\mapsto\varphi\circ f\) is a continuous linear map from \(\mathcal{O}(\mathrm{U};\mathrm{F})\) into \(\mathcal{O}(\mathrm{U};\mathrm{H})\) , denoted \(\varphi_{*}\) .

Let n be a natural number and set \(E = C^{n}\) . Let K be a compact subset of \(C^{n}\) and let U be the decreasing filtered set of open neighborhoods of K. If U, \(U' \in U\) and \(U' \subset U\) , the restriction map of functions from \(\mathcal{O}(U; F)\) into \(\mathcal{O}(U'; F)\) is continuous. The inductive limit of the spaces \(\mathcal{O}(\mathrm{U};\mathrm{F})\) for these maps is denoted \(\mathcal{O}(\mathrm{K};\mathrm{F})\) . The elements of \(\mathcal{O}(\mathrm{K};\mathrm{F})\) are called the germs of holomorphic functions in a neighborhood of K with values in F.

The space \(\mathcal{O}(\mathrm{K};\mathrm{F})\) is endowed with the inductive limit topology of the locally convex topologies of the \(\mathcal{O}(\mathrm{U};\mathrm{F})\) Let X be a locally convex topological vector space and \(\varphi\colon\mathcal{O}(\mathrm{K};\mathrm{F})\to\mathrm{X}\) a map. For every open neighborhood U of K, the map \(\mathcal{O}(\mathrm{U};\mathrm{F})\to\mathrm{X}\) derived from \(\varphi\) by composition with the canonical map \(\mathcal{O}(\mathrm{U};\mathrm{F})\to\mathcal{O}(\mathrm{K};\mathrm{F})\) is denoted \(\varphi^{\mathrm{U}}\) . The map \(\varphi\) is continuous if and only if \(\varphi^{\mathrm{U}}\) is continuous for every U (EVT, II, p.~29, prop. 5, (iii)).

Let m be a natural number. Let L be a compact subset of \(C^{m}\) and let V be an open neighborhood of L. Let \(\varphi\colon V \to C^{n}\) a holomorphic map such that \(\varphi(L) \subset K\) . The continuous linear maps

\[
\mathcal {O} (\mathrm{U}; \mathrm{F}) \xrightarrow {\varphi^ {*}} \mathcal {O} (\stackrel {- 1} {\varphi} (\mathrm{U}); \mathrm{F}) \xrightarrow {\varphi^ {- 1} \stackrel {- 1} {\varphi} (\mathrm{U})} \mathcal {O} (\mathrm{L}; \mathrm{F})
\]

for U an open neighborhood of K, induce a continuous linear map \(\varphi^{*}\colon\mathcal{O}(\mathrm{K};\mathrm{F})\to\mathcal{O}(\mathrm{L};\mathrm{F})\) (loc. cit.).

Let H be a complex Banach space and \(\varphi\colon\mathrm{F}\to\mathrm{H}\) a continuous linear map. The continuous linear maps

\[
\mathcal {O} (\mathrm{U}; \mathrm{F}) \xrightarrow {\varphi_ {*}} \mathcal {O} (\mathrm{U}; \mathrm{H}) \xrightarrow {\varphi^ {\mathrm{U}}} \mathcal {O} (\mathrm{K}; \mathrm{F})
\]

where U ranges over the set of open neighborhoods of K in \(C^{n}\) induce a continuous linear map \(\varphi_{*}\) of \(\mathcal{O}(K;F)\) into \(\mathcal{O}(K;H)\) (loc. cit.). We shall sometimes denote \(\varphi \circ f = \varphi_{*}(f)\) .

For every open neighborhood U of K, the restriction to K is a continuous linear map. \(\mathcal{O}(\mathrm{U};\mathrm{F})\to \mathcal{C}(\mathrm{K};\mathrm{F})\) ; these maps induce a continuous linear map \(\mathcal{O}(\mathrm{K};\mathrm{F})\to \mathcal{C}(\mathrm{K};\mathrm{F})\) called evaluation of germs of holomorphic functions on K.

Let A be a unital complex Banach algebra. The spaces \(\mathcal{O}(\mathrm{U};\mathrm{A})\) and \(\mathcal{O}(\mathrm{K};\mathrm{A})\) are unital algebras. If \(A \neq \{0\}\) one can canonically identify \(\mathcal{O}(\mathrm{U};\mathbf{C})\) (resp. \(\mathcal{O}(\mathrm{K};\mathbf{C})\) ) with the subalgebra \(\mathcal{O}(\mathrm{U};\mathbf{C}) \cdot 1\) of \(\mathcal{O}(\mathrm{U};\mathrm{A})\) (resp. to the subalgebra \(\mathcal{O}(\mathrm{K};\mathbf{C}) \cdot 1\) of \(\mathcal{O}(\mathrm{K};\mathrm{A})\) ). We shall set \(\mathcal{O}(\mathrm{U}) = \mathcal{O}(\mathrm{U};\mathbf{C})\) and \(\mathcal{O}(\mathrm{K}) = \mathcal{O}(\mathrm{K};\mathbf{C})\) .

\subsection{Statement of the main theorem}

Let X be a set. If \(m \leqslant n\) , we shall denote by \(\pi_{m,n}\) the map from \(X^{n}\) into \(X^{m}\) such that \(\pi_{m,n}(\boldsymbol{x}) = (x_{1}, \ldots, x_{m})\) for all \(\boldsymbol{x} = (x_{1}, \ldots, x_{n}) \in \mathrm{X}^{n}\) .

Let A be a unital Banach algebra over C. For every integer \(n \geqslant 1\) and every \(a \in A^{n}\) , we denote by \(\mathrm{Sp}^{n}(\boldsymbol{a})\) the simultaneous spectrum \(\mathrm{Sp}_{\mathrm{A}}^{\{1,\ldots,n\}}(\boldsymbol{a})\) (def. 2 of I, p.~42). It is a compact subset of \(C^{n}\) . For every integer m such that \(1 \leqslant m \leqslant n\) , one has \(\pi_{m,n}(\mathrm{Sp}^{n}(\boldsymbol{a})) = \mathrm{Sp}^{m}(\pi_{m,n}(\boldsymbol{a}))\) (I, p.~41, no. 6). The continuous linear map

\[
\pi_ {m, n} ^ {*} \colon \mathcal {O} (\mathrm{Sp} ^ {m} (\pi_ {m, n} (\boldsymbol {a})); \mathrm{A}) \longrightarrow \mathcal {O} (\mathrm{Sp} ^ {n} (\boldsymbol {a}); \mathrm{A})
\]

is a morphism of unital algebras.

Let A be a commutative unital Banach algebra over C. Let \(n \geqslant 1\) be an integer. One calls holomorphic functional calculus in n variables on A the data, for every \(a \in A^{n}\) , of a map

\[
\Theta_ {\boldsymbol {a}} \colon \mathcal {O} (\mathrm{Sp} ^ {n} (\boldsymbol {a}); \mathrm{A}) \longrightarrow \mathrm{A}
\]

satisfying the conditions:

(CF1) For every \(a \in A^{n}\) , the map \(\Theta_{a}\) is a continuous morphism of unital algebras.

(CF2) If \(\boldsymbol{a}=(a_{1},\ldots,a_{n})\) , and if \(z_{1},\ldots,z_{n}\) denote the germs in a neighborhood of \(\mathrm{Sp}^{n}(\boldsymbol{a})\) of the coordinate functions on \(C^{n}\) , one has

\[
\Theta_ {\boldsymbol {a}} (z _ {1}) = a _ {1}, \dots , \Theta_ {\boldsymbol {a}} (z _ {n}) = a _ {n}.
\]

Remark. --- If the radical of the algebra A is zero, one may omit the continuity condition in (CF1) (cf.~prop. 9 of I, p.~40).

One calls holomorphic functional calculus on A the data, for every integer \(n \geqslant 1\) , of a holomorphic functional calculus in n variables on A, satisfying:

(CF3) Whatever the integers m and n such that \(1 \leqslant m \leqslant n\) , and whatever \(a \in A^{n}\) and \(f \in \mathcal{O}(\mathrm{Sp}^{m}(\pi_{m,n}(a); \mathrm{A})\) , one has

\[
\Theta_ {\boldsymbol {a}} (\pi_ {m, n} ^ {*} (f)) = \Theta_ {\pi_ {m, n} (\boldsymbol {a})} (f).
\]

The object of this paragraph is to prove the following theorem:

THEOREM 1. --- Let A be a commutative complex unital Banach algebra. There exists a unique holomorphic functional calculus on A.

The proof of this theorem will occupy Nos. 3 to 7.

\subsection{Adapted sequences and associated differential forms}

In this number, and up to number 5, we denote by A a commutative complex unital Banach algebra and by n an integer \(\geqslant\) 1.

When we speak of infinitely differentiable functions on an open subset of \(C^{n}\) , it will be a question of functions infinitely differentiable for the underlying structure of real manifold. The notions of differential calculus used will be relative to this structure.

DEFINITION 1. --- Let \(\boldsymbol{a} = (a_{1},\ldots ,a_{n})\in \mathrm{A}^{n}\) , let \(h\colon \mathbf{C}^n\to \mathbf{C}\) be a map and let \(u_{1},\ldots ,u_{n}\) be maps from \(\mathbf{C}^n\) into A. One says that the sequence \((h,u_1,\dots ,u_n)\) is adapted to \(\boldsymbol{a}\) if

\begin{enumerate}
\def\labelenumi{(\roman{enumi})}
\item
  The map h is infinitely differentiable, with compact support, and equal to 1 in a neighborhood of \(\mathrm{Sp}^n (\boldsymbol {a})\) ;
\item
  The maps \(u_1, \ldots, u_n\) are infinitely differentiable;
\item
  For every \(\boldsymbol{z} = (z_{1},\dots ,z_{n})\in \mathbf{C}^{n}\) , we have
\end{enumerate}

\[
h (\boldsymbol {z}) + (z _ {1} - a _ {1}) u _ {1} (\boldsymbol {z}) + \dots + (z _ {n} - a _ {n}) u _ {n} (\boldsymbol {z}) = 1.\tag{1}
\]

The differential form of degree 2n on \(C^{n}\) , with coefficients in A, defined by

\[
\omega = \bigwedge_ {i = 1} ^ {n} (d u _ {i} \wedge d z _ {i})
\]

is called the differential form associated with \((h, u_1, \ldots, u_n)\) .

If \((h, u_{1}, \ldots, u_{n})\) is adapted to a, then by differentiating (1) we obtain the equality

\[
d h = - \sum_ {i = 1} ^ {n} u _ {i} d z _ {i} - \sum_ {i = 1} ^ {n} (z _ {i} - a _ {i}) d u _ {i}\tag{2}
\]

whence, for every i such that \(1 \leqslant i \leqslant n\) , the relation

\[
d h \wedge d z _ {i} \wedge \bigwedge_ {j \neq i} (d u _ {j} \wedge d z _ {j}) = - (z _ {i} - a _ {i}) \omega .\tag{3}
\]

Lemma 1. --- Let U be an open subset of \(\mathbf{C}^{n}\) Let K be a compact subset of U. There exists an infinitely differentiable function h from \(\mathbf{C}^{n}\) in C, equal to 1 on K and with compact support contained in U.

Let V be a relatively compact open neighborhood of K such that \(\overline{V}\) is contained in U (TG, I, p.~65, prop. 10). There exists an infinitely differentiable function h of \(C^{n}\) in C whose support is contained in V and which equals 1 on K (VAR, R1, p.~40, 5.3.6). This function has the required properties.

Example. --- Suppose that n = 1. Let \(a \in A\) . For every open neighborhood U of \(\mathrm{Sp}(a)\) , there exists an infinitely differentiable function h from C into C with compact support contained in U, equal to 1 in a neighborhood of \(\mathrm{Sp}(a)\) (VAR, R1, p.~40, 5.3.6). Set

\[
u (z) = (1 - h (z)) (z - a) ^ {- 1}
\]

for \(z \in \mathbf{C} - \mathrm{Sp}(a)\) and \(u(z) = 0\) if \(z \in \mathrm{Sp}(a)\) . The pair \((h, u)\) is adapted to a and the associated differential form is \(\omega = du \wedge dz\) .

Lemma 2. --- Let \(\boldsymbol{a} = (a_{1},\ldots ,a_{n})\in \mathrm{A}^{n}\) There exist infinitely differentiable maps \(v_{1},\ldots ,v_{n}\) of \(\mathbf{C}^n -\mathrm{Sp}^n (\boldsymbol {a})\) in A such that

\[
(z _ {1} - a _ {1}) v _ {1} (\boldsymbol {z}) + \dots + (z _ {n} - a _ {n}) v _ {n} (\boldsymbol {z}) = 1
\]

for all \(z = (z_{1}, \ldots, z_{n}) \in \mathbf{C}^{n} - \mathrm{Sp}^{n}(\boldsymbol{a})\) .

Let \(\boldsymbol{w} = (w_{1}, \ldots, w_{n}) \in \mathbf{C}^{n} - \operatorname{Sp}^{n}(\boldsymbol{a})\) . By definition of the simultaneous spectrum, there exist \(b_{1}, \ldots, b_{n}\) in A such that

\[
(w _ {1} - a _ {1}) b _ {1} + \dots + (w _ {n} - a _ {n}) b _ {n} = 1
\]

(I, p.~41, no. 6). There exists an open neighborhood \(W_{w}\) of w such that the element \((z_{1}-a_{1})b_{1}+\cdots+(z_{n}-a_{n})b_{n}\) of A is invertible if \(z=(z_{1},\ldots,z_{n})\) belongs to \(W_{w}\) . For every integer j such that \(1\leqslant j\leqslant n\) and every z in \(W_{w}\) , then let

\[
u _ {j} (\boldsymbol {z}) = b _ {j} \left(\sum_ {i = 1} ^ {n} (z _ {i} - a _ {i}) b _ {i}\right) ^ {- 1}.
\]

The functions \(u_{1}, u_{2}, \ldots, u_{n}\) of \(W_{w}\) in A thus defined are infinitely differentiable in \(W_{w}\) , and we have

\[
(z _ {1} - a _ {1}) u _ {1} (\boldsymbol {z}) + \dots + (z _ {n} - a _ {n}) u _ {n} (\boldsymbol {z}) = 1
\]

for all z in \(W_{w}\) .

Since the family \((\mathrm{W}_{\boldsymbol{w}})_{\boldsymbol{w}\in\mathbf{C}^{n}-\mathrm{Sp}^{n}(\boldsymbol{a})}\) is an open covering of \(\mathbf{C}^{n}-\mathrm{Sp}^{n}(\boldsymbol{a})\) , there exists a locally finite open refinement \(\mathscr{W}=(\mathrm{W}_{\lambda})_{\lambda\in\mathrm{L}}\) (TG, I, p.~70, th. 5) and, for every \(\lambda\in L\) , functions \(u_{1\lambda},\ldots,u_{n\lambda}\) , with values in A, defined and infinitely differentiable in \(W_{\lambda}\) , such that \((z_{1}-a_{1})u_{1\lambda}(\boldsymbol{z})+\cdots+(z_{n}-a_{n})u_{n\lambda}(\boldsymbol{z})=1\) for all \(\boldsymbol{z}\) into \(W_{\lambda}\) . Let \((f_{\lambda})_{\lambda \in L}\) a partition of unity subordinate to the covering \(\mathscr{W}\) consisting of infinitely differentiable functions (VAR, R1, p.~40, 5.3.6). Let \(i\) be an integer such that \(1 \leqslant i \leqslant n\) . For every \(\lambda \in L\) , let \(u_{i\lambda}'\) the map from \(\mathbf{C}^n - \mathrm{Sp}^n(\boldsymbol{a})\) in A obtained by extending by 0 the function \(f_{\lambda}u_{i\lambda}\) into \((\mathbf{C}^n - \mathrm{Sp}^n(\boldsymbol{a})) - W_{\lambda}\) . The functions \(u_{i\lambda}'\) are infinitely differentiable. Since the family \((\mathrm{Supp}(u_{i\lambda}')_{\lambda \in L}\) is locally finite, the function \(v_i = \sum_{\lambda \in L} u_{i\lambda}'\) is defined and infinitely differentiable in \(\mathbf{C}^n - \mathrm{Sp}^n(\boldsymbol{a})\) .

Let \(z \in \mathbf{C}^n - \mathrm{Sp}^n(\boldsymbol{a})\) . Let L' denote the finite set of \(\lambda \in \mathrm{L}\) such that \(z \in \mathrm{W}_{\lambda}\) . Then

\[
\begin{array}{c} \sum_ {i = 1} ^ {n} (z _ {i} - a _ {i}) v _ {i} (\boldsymbol {z}) = \sum_ {\lambda \in \mathrm{L} ^ {\prime}} \sum_ {i = 1} ^ {n} (z _ {i} - a _ {i}) u _ {i \lambda} ^ {\prime} (\boldsymbol {z}) \\ = \sum_ {\lambda \in \mathrm{L} ^ {\prime}} f _ {\lambda} (\boldsymbol {z}) \sum_ {i = 1} ^ {n} (z _ {i} - a _ {i}) u _ {i \lambda} (\boldsymbol {z}) = \bigg (\sum_ {\lambda \in \mathrm{L} ^ {\prime}} f _ {\lambda} (\boldsymbol {z}) \bigg) \cdot 1 = 1. \end{array}
\]

Lemma 3. --- Let \(\boldsymbol{a} = (a_{1},\ldots ,a_{n})\in \mathrm{A}^{n}\) . Let \(h\) a map of \(\mathbf{C}^n\) into \(\mathbf{C}\) infinitely differentiable, equal to 1 in a neighborhood of \(\mathrm{Sp}^n (\boldsymbol {a})\) and with compact support. There exist infinitely differentiable maps \(u_{1},\ldots ,u_{n}\) of \(\mathbf{C}^n\) in A such that the sequence \((h,u_1,\dots,u_n)\) is adapted to \(\boldsymbol{a}\) .

Let \(v_{1},\ldots,v_{n}\) be infinitely differentiable maps from \(\mathbf{C}^{n}-\mathrm{Sp}^{n}(\boldsymbol{a})\) into A such that

\[
\sum_ {j = 1} ^ {n} (z _ {j} - a _ {j}) v _ {j} (\boldsymbol {z}) = 1
\]

for \(z\) into \(\mathbf{C}^n - \mathrm{Sp}^n(\boldsymbol{a})\) (lemma 2). Let \(i\) be an integer such that \(1 \leqslant i \leqslant n\) . Set \(u_i(z) = (1 - h(z))v_i(z)\) if \(z \in \mathbf{C}^n - \mathrm{Sp}^n(\boldsymbol{a})\) and \(u_i(z) = 0\) if \(z \in \mathrm{Sp}^n(\boldsymbol{a})\) The mappings \(u_i\) are infinitely differentiable in \(\mathbf{C}^n - \mathrm{Sp}^n(\boldsymbol{a})\) and vanish in a neighborhood of \(\mathrm{Sp}^n(\boldsymbol{a})\) , hence infinitely differentiable in \(\mathbf{C}^n\) . Equality (1) holds in \(\mathrm{Sp}^n(\boldsymbol{a})\) because the functions \(u_i\) are zero on \(\mathrm{Sp}^n(\boldsymbol{a})\) and \(h\) is equal to 1 in a neighborhood of \(\mathrm{Sp}^n(\boldsymbol{a})\) It is also true on \(\mathbf{C}^n - \mathrm{Sp}^n(\boldsymbol{a})\) by construction.

Lemma 4. — Let \(a \in A^{n}\) . Let \((h, u_{1}, \ldots, u_{n})\) be a sequence adapted to a and \(\omega\) the associated differential form.

\begin{enumerate}
\def\labelenumi{\alph{enumi})}
\tightlist
\item
  For \(i = 1,2,\ldots,n\) , there exists a differential form \(\beta_{i}\) on \(\mathbf{C}^{n}\) , of degree \(n - 1\) , with coefficients in A, such that
\end{enumerate}

\[
(z _ {i} - a _ {i}) \omega = d (h \beta_ {i} \wedge d z _ {1} \wedge \dots \wedge d z _ {n});
\]

\begin{enumerate}
\def\labelenumi{\alph{enumi})}
\setcounter{enumi}{1}
\item
  The differential form \(\omega\) has compact support contained in the support of \(h\) ;
\item
  There exists a differential form \(\beta\) on \(\mathbf{C}^{n}\) , of degree \(n - 1\) , with coefficients in A, such that
\end{enumerate}

\[
(n + 1) h \omega - \omega = d (h \beta \wedge d z _ {1} \wedge \dots \wedge d z _ {n}).
\]

Let \(i\) be an integer such that \(1 \leqslant i \leqslant n\) . There exists \(\varepsilon_i \in \{-1, 1\}\) such that

\[
\varepsilon_ {i} \bigwedge_ {j \neq i} d u _ {j} \wedge d z _ {1} \wedge \dots \wedge d z _ {n} = d z _ {i} \wedge \bigwedge_ {j \neq i} (d u _ {j} \wedge d z _ {j}).
\]

Let us set \(\beta_{i} = \varepsilon_{i} \bigwedge_{j \neq i} du_{j}\) , so that the left-hand term in this formula is \(\beta_{i} \wedge dz_{1} \wedge \cdots \wedge dz_{n}\) and that \(d\beta_{i} = 0\) . Thus

\[
d \left(h \beta_ {i} \wedge d z _ {1} \wedge \dots \wedge d z _ {n}\right) = d h \wedge \beta_ {i} \wedge d z _ {1} \wedge \dots \wedge d z _ {n} =
\]

\[
d h \wedge d z _ {i} \wedge \bigwedge_ {j \neq i} (d u _ {j} \wedge d z _ {j}) = (z _ {i} - a _ {i}) \omega ,
\]

by formula (3), whence the assertion \(a\) ).

From assertion \(a\) ) and formula (1) we deduce the relation

\[
\begin{array}{l} \omega = h \omega + (1 - h) \omega = h \omega + \sum_ {i = 1} ^ {n} (z _ {i} - a _ {i}) u _ {i} \omega \\ \qquad \qquad \qquad = h \omega + \sum_ {i = 1} ^ {n} u _ {i}   d (h \beta_ {i} \wedge d z _ {1} \wedge \dots \wedge d z _ {n}), \end{array}
\]

whence \(\operatorname{Supp}(\omega) \subset \operatorname{Supp}(h)\) which proves \(b\) ).

Finally, set

\[
\beta = \sum_ {i = 1} ^ {n} \varepsilon_ {i} u _ {i} \bigwedge_ {j \neq i} d u _ {j} = \sum_ {i = 1} ^ {n} u _ {i} \beta_ {i}, \text {   and   } \tau = h \beta d z _ {1} \wedge \dots \wedge d z _ {n}.
\]

We have \(d\beta = \sum_{i}du_{i}\wedge \beta_{i}\) and therefore

\[
d \beta \wedge d z _ {1} \wedge \dots \wedge d z _ {n} = \sum_ {i} d u _ {i} \wedge d z _ {i} \wedge \bigwedge_ {j \neq i} (d u _ {j} \wedge d z _ {j}) = n \omega .
\]

Thus

\[
\begin{array}{l} d \tau = d h \wedge \beta \wedge d z _ {1} \wedge \dots \wedge d z _ {n} + h d \beta \wedge d z _ {1} \wedge \dots \wedge d z _ {n} \\ = \sum_ {i = 1} ^ {n} u _ {i} d h \wedge d z _ {i} \wedge \bigwedge_ {j \neq i} (d u _ {j} \wedge d z _ {j}) + n h \omega \\ = - \sum_ {i = 1} ^ {n} u _ {i} (z _ {i} - a _ {i}) \omega + n h \omega = (h - 1) \omega + n h \omega = (n + 1) h \omega - \omega , \end{array}
\]

Given formulas (3) and (1), whence c).

We now propose to study how the differential form \(\omega\) associated with a sequence adapted to \(\boldsymbol{a}\) varies as a function of this sequence. We will say that sequences \((h, u_1, \ldots, u_n)\) and \((h', u_1', \ldots, u_n')\) adapted to \(\boldsymbol{a}\) are linked if there exists a differential form \(\psi\) of degree \(n-1\) on \(\mathbf{C}^n\) , with coefficients in A and with support contained in the union of the supports of \(h\) and of \(h'\) , such that the associated differential forms \(\omega\) and \(\omega'\) satisfy

\[
\omega - \omega^ {\prime} = d (\psi \wedge d z _ {1} \wedge d z _ {2} \wedge \dots \wedge d z _ {n}).
\]

Let us begin with an elementary modification:

Lemma 5. --- Let \(\boldsymbol{a} \in \mathrm{A}^{n}\) , let \((h, u_{1}, \ldots, u_{n})\) be a sequence adapted to \(\boldsymbol{a}\) , and let \(\omega\) the associated differential form.

Let \(w\) an infinitely differentiable map from \(\mathbf{C}^n\) into \(A\) and let \(i\) and \(j\) be distinct integers between 1 and \(n\) . Define \(u_1', \ldots, u_n'\) by

\[
\begin{array}{c} u _ {i} ^ {\prime} = u _ {i} + (z _ {j} - a _ {j}) w, \quad u _ {j} ^ {\prime} = u _ {j} - (z _ {i} - a _ {i}) w, \\ u _ {k} ^ {\prime} = u _ {k} \text {for} k \neq i, j. \end{array}
\]

Then the sequence \((h,u_{1}^{\prime},\ldots,u_{n}^{\prime})\) is adapted to a and is linked to the sequence \((h,u_{1},\ldots,u_{n})\) .

We denote \(dz = dz_{1} \wedge \cdots \wedge dz_{n}\) . Since

\[
\begin{array}{c} \sum_ {k = 1} ^ {n} (z _ {k} - a _ {k}) u _ {k} ^ {\prime} (\boldsymbol {z}) = \sum_ {k = 1} ^ {n} (z _ {k} - a _ {k}) u _ {k} (\boldsymbol {z}) + w (\boldsymbol {z}) (z _ {j} - a _ {j}) (z _ {i} - a _ {i}) \\ - w (\boldsymbol {z}) (z _ {i} - a _ {i}) (z _ {j} - a _ {j}) = 1 - h (\boldsymbol {z}) \end{array}
\]

for all \(z \in C^{n}\) the sequence \((h, u_{1}', \ldots, u_{n}')\) is adapted to a. Moreover, we have

\[
\begin{array}{l} d u _ {i} ^ {\prime} \wedge d u _ {j} ^ {\prime} \wedge d z _ {1} \wedge \dots \wedge d z _ {n} = \\ \left(d u _ {i} + w d z _ {j} + (z _ {j} - a _ {j}) d w\right) \wedge \left(d u _ {j} - w d z _ {i} - (z _ {i} - a _ {i}) d w\right) \wedge d \boldsymbol {z} \\ = \left(d u _ {i} \wedge d u _ {j} - (z _ {i} - a _ {i}) d u _ {i} \wedge d w - (z _ {j} - a _ {j}) d u _ {j} \wedge d w\right) \wedge d \boldsymbol {z} \end{array}
\]

There therefore exists \(\varepsilon\in\{-1,1\}\) such that \(\varepsilon(\omega-\omega')\) is equal to

\[
\begin{array}{l} d u _ {i} ^ {\prime} \wedge d u _ {j} ^ {\prime} \wedge \bigwedge_ {k \neq i, j} d u _ {k} ^ {\prime} \wedge d \boldsymbol {z} - d u _ {i} \wedge d u _ {j} \wedge \bigwedge_ {k \neq i, j} d u _ {k} \wedge d \boldsymbol {z} \\ = - \Big ((z _ {i} - a _ {i}) d u _ {i} \wedge d w + (z _ {j} - a _ {j}) d u _ {j} \wedge d w \Big) \wedge \bigwedge_ {k \neq i, j} d u _ {k} \wedge d \boldsymbol {z} \\ = - \Big (\sum_ {k = 1} ^ {n} (z _ {k} - a _ {k})   d u _ {k} \Big) \wedge d w \wedge \bigwedge_ {k \neq i, j} d u _ {k} \wedge d \boldsymbol {z} \end{array}
\]

and, taking (2) into account, this is equal to

\[
d h \wedge d w \wedge \bigwedge_ {k \neq i, j} d u _ {k} \wedge d \boldsymbol {z} = d \left(h d w \wedge \bigwedge_ {k \neq i, j} d u _ {k} \wedge d \boldsymbol {z}\right),
\]

whence the result.

Lemma 6. — Let \(\mathbf{a} \in \mathrm{A}^n\) . All sequences adapted to \(\mathbf{a}\) are linked.

Let \((h,u_{1},\ldots,u_{n})\) and \((h',u_{1}',\ldots,u_{n}')\) of sequences adapted to a, and denote \(\omega\) and \(\omega'\) the associated differential forms.

Let us define the infinitely differentiable maps

\[
\begin{array}{c} {w _ {i j} = u _ {i} ^ {\prime} u _ {j} - u _ {i} u _ {j} ^ {\prime},} \\ {s _ {i} = u _ {i} ^ {\prime} h - u _ {i} h ^ {\prime},} \end{array}
\]

\[
\begin{array}{l} 1 \leqslant i \leqslant n, 1 \leqslant j \leqslant n, \\ 1 \leqslant i \leqslant n, \end{array}
\]

so that \(w_{ji} = -w_{ij}\) , and \(\operatorname{Supp}(s_i) \subset \operatorname{Supp}(h) \cup \operatorname{Supp}(h')\) .

Let us set \(u_i'' = u_i' - s_i\) , \(\boldsymbol{u} = (u_1, \ldots, u_n)\) and \(\boldsymbol{u}'' = (u_1'', \ldots, u_n'')\) Let us also note \(\boldsymbol{v}_{ij}\) the map from \(\mathbf{C}^n\) into \(A^n\) whose \(i\) -th component is \((z_j - a_j) w_{ij}\) , whose \(j\) -th component is \((z_i - a_i) w_{ji} = -(z_i - a_i) w_{ij}\) and whose other components are zero. Then we have

\[
\boldsymbol {u} ^ {\prime \prime} = \boldsymbol {u} + \sum_ {i <   j} \boldsymbol {v} _ {i j}.
\]

Indeed, for every integer k such that \(1 \leqslant k \leqslant n\) , the k-th component of the right-hand side is

\[
\begin{array}{c} u _ {k} + \sum_ {i = 1} ^ {k - 1} (z _ {i} - a _ {i}) w _ {k i} + \sum_ {j = k + 1} ^ {n} (z _ {j} - a _ {j}) w _ {k j} = u _ {k} + \sum_ {i = 1} ^ {n} (z _ {i} - a _ {i}) w _ {k i} \\ = u _ {k} + u _ {k} ^ {\prime} \sum_ {i = 1} ^ {n} (z _ {i} - a _ {i}) u _ {i} - u _ {k} \sum_ {i = 1} ^ {n} (z _ {i} - a _ {i}) u _ {i} ^ {\prime} \\ = u _ {k} + (1 - h) u _ {k} ^ {\prime} - (1 - h ^ {\prime}) u _ {k} = u _ {k} ^ {\prime} - s _ {k}. \end{array}
\]

By induction, we deduce from Lemma 5, applied to the integers \(i\) and \(j\) and to the maps \(w_{ij}\) , that the sequence \((h,u_1'',\ldots,u_n'')\) is adapted to \(\boldsymbol{a}\) and is linked to \((h,u_1,\ldots,u_n)\) . Let \(\omega''\) the differential form associated with \((h,u_1'',\ldots,u_n'')\) . Since \(u_i'' = u_i' - s_i\) , one has

\[
\begin{array}{r} \omega^ {\prime \prime} - \omega^ {\prime} = d (u _ {1} ^ {\prime} - s _ {1}) \wedge d z _ {1} \wedge \dots \wedge d (u _ {n} ^ {\prime} - s _ {n}) \wedge d z _ {n} - \\ d u _ {1} ^ {\prime} \wedge d z _ {1} \wedge \dots \wedge d u _ {n} ^ {\prime} \wedge d z _ {n}, \end{array}
\]

which is expressed as a linear combination, with coefficients 1 or -1, of differential forms of the form

\[
\xi_ {\mathrm{I} _ {1}, \mathrm{I} _ {2}} = \bigwedge_ {i \in \mathrm{I} _ {1}} d s _ {i} \wedge \bigwedge_ {i \in \mathrm{I} _ {2}} d u _ {i} ^ {\prime} \wedge d z _ {1} \wedge \dots \wedge d z _ {n},
\]

where \(I_{1}\) (resp. \(I_{2}\) ) is a nonempty subset (resp. a subset) of \(\{1,\ldots,n\}\) . Each differential form \(\xi_{I_{1},I_{2}}\) can also be written in the form

\[
d (\widetilde {\psi} \wedge d z _ {1} \wedge \dots \wedge d z _ {n}),
\]

where the support of the differential form \(\widetilde{\psi}\) is contained in the support of \(s_i\) for all \(i \in \mathrm{I}_1\) . Since \(\mathrm{I}_1\) is nonempty, this support is contained in \(\operatorname{Supp}(h) \cup \operatorname{Supp}(h')\) . Consequently, \((h, u_1'', \ldots, u_n'')\) is linked to \((h', u_1', \ldots, u_n')\) , and the lemma follows by writing

\[
\omega - \omega^ {\prime} = (\omega - \omega^ {\prime \prime}) + (\omega^ {\prime \prime} - \omega^ {\prime}).
\]

\subsection{\texorpdfstring{Construction of the maps \(\Theta_{a}\)}{Construction of the maps \textbackslash Theta\_\{a\}}}

Let \(\boldsymbol{a}=(a_{1},\ldots,a_{n})\in\mathrm{A}^{n}\) and let U be an open neighborhood of \(\mathrm{Sp}^{n}(\boldsymbol{a})\) Let h be an infinitely differentiable function, equal to 1 in a neighborhood of \(\mathrm{Sp}^{n}(\boldsymbol{a})\) and such that the support of h is compact and contained in U (Lemma 1 of I, p.~52). By Lemma 3 of I, p.~54, there exist infinitely differentiable functions \((u_{1},\ldots,u_{n})\) of \(C^{n}\) in A such that the sequence \((h,u_{1},\ldots,u_{n})\) adapted to a. Let \(\omega\) be the associated differential form; it has compact support contained in U (Lemma 4 of I, p.~54). There exists an infinitely differentiable function \(\psi\) with compact support in U and with values in A such that

\[
\omega = \psi d x _ {1} \wedge d y _ {1} \wedge \dots \wedge d x _ {n} \wedge d y _ {n}
\]

where \(x_{j} + iy_{j}\) are the coordinate functions on \(C^{n}\) identified with \(R^{2n}\) . Let \(\mu\) the Lebesgue measure on \(R^{2n}\) .

Let \(f \in \mathcal{C}(\mathrm{U};\mathrm{A})\) The differential form \(f\omega|\mathrm{U}\) on \(\mathrm{U}\) is continuous and compactly supported. The vector measure associated with this differential form (VAR, R2, 10.4.3 and 10.4.4) is the vector measure \(f\psi \cdot \mu\) ; it is a basic measure \(\mu\) (INT, VI, §2, no. 4, def. 4), its support is compact and contained in the support of \(\omega\) This measure is bounded relative to the norm of A (INT, VI, §2, no. 4, prop. 8); we denote by \(\| f\omega\|\) the positive measure on \(\mathbf{C}^n\) associated with it (INT, VI, §2, no. 3, def. 3).

By INT, VI, §2, no. 4, prop. 8, b), we have

\[
\| f \omega \| = \| f \psi \cdot \mu \| = \| f \psi \| \cdot \mu = \| f \| \| \omega \|.
\]

In particular, the integral of the differential form \(f\omega\) over U satisfies

\[
\left\| \int_ {\mathrm{U}} f \omega \right\| \leqslant \int_ {\mathrm{U}} \| f \| \| \omega \|
\]

(INT, VI, §2, no. 3, prop. 5).

Lemma 7. --- For every function \(f \in \mathcal{O}(\mathrm{U};\mathrm{A})\) , the integral \(\int_{\mathrm{U}} f\omega\) is an element of A which depends only on \(\boldsymbol{a}\) and on the germ of \(f\) in a neighborhood of \(\mathrm{Sp}^{n}(\boldsymbol{a})\) . It satisfies the inequality

\[
\left\| \int_ {\mathrm{U}} f \omega \right\| \leqslant \left(\int_ {\mathrm{U}} \| \omega \|\right) \sup _ {z \in \operatorname{Supp} (h)} \| f (z) \|,\tag{4}
\]

and

\[
\int_ {\mathrm{U}} (a f) \omega = a \int_ {\mathrm{U}} f \omega
\]

for all \(a \in \mathbf{A}\) .

We saw above that the integral is defined for \(f \in \mathcal{C}(\mathrm{U};\mathrm{A})\) and satisfies

\[
\left\| \int_ {\mathrm{U}} f \omega \right\| \leqslant \int_ {\mathrm{U}} \| f \| \| \omega \| \leqslant \left(\int_ {\mathrm{U}} \| \omega \|\right) \sup _ {z \in \operatorname{Supp} (h)} \| f (z) \|.
\]

Moreover, for every \(a \in \mathrm{A}\) and every \(f \in \mathcal{C}(\mathrm{U};\mathrm{A})\) , one has

\[
\int_ {\mathrm{U}} (a f) \omega = a \int_ {\mathrm{U}} f \omega
\]

(INT, VI, §2, no. 2, prop. 2 applied to multiplication by a).

Let \((h', u_{1}', \ldots, u_{n}')\) a sequence adapted to a such that \(\operatorname{Supp}(h') \subset \operatorname{U}\) . Let \(\omega'\) the associated differential form. By lemma 6 of I, p.~57, there exists a differential form \(\psi\) on \(C^{n}\) of degree n-1, with coefficients in A and with support contained in \(\operatorname{Supp}(h) \cup \operatorname{Supp}(h') \subset \operatorname{U}\) , such that

\[
\omega - \omega^ {\prime} = d (\psi \wedge d z _ {1} \wedge \dots \wedge d z _ {n}).
\]

Let \(f \in \mathcal{O}(\mathrm{U};\mathrm{A})\) . The map \(f\) being holomorphic, we have

\[
d f = \sum_ {i = 1} ^ {n} \frac {\partial f}{\partial z _ {i}} d z _ {i},
\]

(VAR, R2, p.~24, 8.8.9) and hence

\[
f \left(\omega - \omega^ {\prime}\right) = f d \left(\psi \wedge d z _ {1} \wedge \dots \wedge d z _ {n}\right) = d \left(f \psi \wedge d z _ {1} \wedge \dots \wedge d z _ {n}\right).
\]

By Stokes' formula (VAR, R2, p.~48, 11.2.3), we then have

\[
\int_ {\mathrm{U}} f (\omega - \omega^ {\prime}) = 0.
\]

Thus the element \(\int_{\mathrm{U}} f\omega\) does not depend on the choice of the sequence \((h, u_1, \ldots, u_n)\) . To conclude, let us show that \(\int_{\mathrm{U}} f\omega\) depends only on the germ of \(f\) in a neighborhood of \(\mathrm{Sp}^n(\boldsymbol{a})\) . Let U and \(\mathrm{U}'\) be open neighborhoods of \(\mathrm{Sp}^n(\boldsymbol{a})\) . Let \(f \in \mathcal{O}(\mathrm{U}; \mathrm{A})\) and \(f' \in \mathcal{O}(\mathrm{U}', \mathrm{A})\) such that \(f\) and \(f'\) coincide on an open neighborhood \(\mathrm{U}''\) of \(\mathrm{Sp}^n(\boldsymbol{a})\) . There exists a map \(h\) of \(\mathbf{C}^n\) into \(\mathbf{C}\) infinitely differentiable, equal to 1 in a neighborhood of \(\mathrm{Sp}^n(\boldsymbol{a})\) , with compact support contained in \(\mathrm{U}''\) (lemma 1 of I, p.~52), and there exists \((u_1, \ldots, u_n)\) such that the sequence \((h, u_1, \ldots, u_n)\) is adapted to \(\boldsymbol{a}\) (lemma 3 of I, p.~54). Let \(\omega\) be the associated differential form. Since \(\mathrm{Supp}(\omega) \subset \mathrm{Supp}(h) \subset \mathrm{U}''\) (lemma 4, b) of I, p.~54), we have

\[
\int_ {\mathrm{U}} f \omega = \int_ {\mathrm {U^ {\prime\prime}}} f \omega = \int_ {\mathrm {U^ {\prime\prime}}} f ^ {\prime} \omega = \int_ {\mathrm {U^ {\prime}}} f ^ {\prime} \omega ,
\]

which completes the proof.

This lemma shows that there exists a unique A-linear map \(\Theta_{\boldsymbol{a}}\) of \(\mathcal{O}(\mathrm{Sp}^{n}(\boldsymbol{a});\mathrm{A})\) to A such that

\[
\Theta_ {\boldsymbol {a}} (f) = \frac {n !}{(2 i \pi) ^ {n}} \int_ {\mathrm{U}} \widetilde {f} \omega\tag{5}
\]

for every open set U and every representative \(\tilde{f} \in \mathcal{O}(\mathrm{U};\mathrm{A})\) of a germ \(f \in \mathcal{O}(\mathrm{Sp}^n (\boldsymbol {a});\mathrm{A})\) . The linear map \(\Theta_{\boldsymbol{a}}\) is continuous by inequality (4) and EVT, II, p.~29, prop. 5.
